\documentclass[12pt]{article}
\usepackage{amsmath}
\usepackage{amsfonts}
\usepackage{float}
\usepackage{fancyhdr}
\usepackage{graphicx}
\usepackage{setspace}
\usepackage[colorlinks=true,linkcolor=blue, citecolor=red]{hyperref}
\usepackage{url}
\usepackage[top=.75in, left=.75in, right=.75in, bottom=1in]{geometry}

% For algorithm
\usepackage{algorithm}
\usepackage{algpseudocode}

% ============ CODE ============
\usepackage{listings}
\usepackage{xcolor}
\definecolor{codegreen}{rgb}{0,0.6,0}
\definecolor{codegray}{rgb}{0.5,0.5,0.5}
\definecolor{codepurple}{rgb}{0.58,0,0.82}
\definecolor{backcolour}{rgb}{0.95,0.95,0.92}

% Styling for the code.
\lstdefinestyle{mystyle}{
    backgroundcolor=\color{backcolour},   
    commentstyle=\color{codegreen},
    keywordstyle=\color{magenta},
    numberstyle=\tiny\color{codegray},
    stringstyle=\color{codepurple},
    basicstyle=\ttfamily\footnotesize,
    breakatwhitespace=false,         
    breaklines=true,                 
    captionpos=b,                    
    keepspaces=true,                 
    numbers=left,                    
    numbersep=5pt,                  
    showspaces=false,                
    showstringspaces=false,
    showtabs=false,                  
    tabsize=2
}
\lstset{style=mystyle}

% Table packages and Unicode / Font support
\usepackage{array}
\usepackage{longtable}
\usepackage{enumitem}
\usepackage{booktabs}
\usepackage[utf8]{inputenc}
\DeclareUnicodeCharacter{2264}{\ensuremath{\le}}
\DeclareUnicodeCharacter{2194}{\ensuremath{\leftrightarrow}}
\usepackage[T1]{fontenc}
\usepackage{lmodern}
\graphicspath{{../../../../templates/}{../../prototype/}{figures/}{./}}
\setlength{\emergencystretch}{3em}
\setlist[enumerate]{label=\arabic*.,leftmargin=*}
\setlist[itemize]{leftmargin=*}

% ============ HEADER AND FOOTER ============
\setlength{\headheight}{29.43912pt}

\pagestyle{fancy}
\fancyfoot{}
\fancyfoot[R]{Page \thepage}

\setlength{\parindent}{0pt}
\renewcommand{\baselinestretch}{1.5}

\newcommand{\coursename}{Introduction of Software Engineering}
\newcommand{\reportname}{Studitopia}
\newcommand{\reporttitle}{\href{https://github.com/RoundFloor/Studitopia.git}{Studitopia}}

\newcommand{\studentname}{%
Nguyen Thanh An (24125024)\\
Nguyen Chi Bao (24125025)\\
Tran Nguyen Khanh Binh (24125026)\\
Tran Canh Anh Tuan (24125048)\\
Nguyen Le Quoc Huy (24125100)
}
\newcommand{\teachername}{Ho Tuan Thanh, Nguyen Thi Minh Tuyen}

% Header
\lhead{\reporttitle}
\rhead{
University of Science - VNUHCM\\
\coursename
}

% Footer
\newcommand{\leftfooter}{}
\lfoot{\leftfooter}

% ============ DOCUMENT ============
\begin{document}

\pagenumbering{roman}
\begin{titlepage}
\newcommand{\HRule}{\rule{\linewidth}{0.5mm}}
\centering

\textsc{\LARGE VNU-HCM}\\[0.5cm]
\textsc{\Large University of Science}\\[0.5cm]
\textsc{\large Faculty of Information Technology}\\[0.5cm]
\textsc{Computer Science}\\[0.5cm]

\HRule \\[0.4cm]
{ 
\huge{\bfseries{\reporttitle}}\\[0.5cm]
\large{\bfseries{Learning Resource Library --- Use Case Diagram}}
}\\[0.4cm]
\HRule \\[0.5cm]

\textbf{\large Subject: \coursename}\\[0.5cm]

\begin{minipage}[t]{0.55\textwidth}
\begin{flushleft} \large
\emph{Students:}\\
\studentname
\end{flushleft}
\end{minipage}
~
\begin{minipage}[t]{0.4\textwidth}
\begin{flushright} \large
\emph{Supervising Teachers:} \\
\teachername
\end{flushright}
\end{minipage}\\[1cm]

{\large \today}\\[1cm]

\includegraphics[scale=.20]{img/hcmus-logo.png}\\[1cm] 

\vfill
\end{titlepage}	

\tableofcontents
\pagebreak

\pagenumbering{arabic}
\setcounter{page}{1}

\hypertarget{sec:learning-resource-library-use-case-diagram}{}
\addcontentsline{toc}{section}{Learning Resource Library --- Use Case Diagram}
\section*{Learning Resource Library --- Use Case Diagram}
\par\medskip

\textbf{Version:} \texttt{1.0}
\par\medskip

\textbf{Performed by:} [Nguyen Le Quoc Huy] | \textbf{Edited by:} [Nguyen Le Quoc Huy]
\par\medskip

Merged from the focused use case diagrams in this functional group.
\par\medskip

\begin{center}
\includegraphics[width=0.85\textwidth,height=0.75\textheight,keepaspectratio]{"figures/diagram_1.png"}
\end{center}
\par\medskip

\textbf{Diagram notation.} Person shape = human actor; blue = base use case; purple = included; green = extending; solid line = association; dashed arrow = include or extend.
\par\medskip

\pagebreak
\hypertarget{sec:uc-06-share-learning-resource}{}
\addcontentsline{toc}{section}{UC-06 --- Share Learning Resource}
\section*{UC-06 --- Share Learning Resource}
\par\medskip

\textbf{Version:} \texttt{2.0}
\par\medskip

\textbf{Performed by:} [Tran Nguyen Khanh Binh] | \textbf{Edited by:} [Nguyen Le Quoc Huy]
\par\medskip

\begingroup
\small
\renewcommand{\baselinestretch}{1.0}\selectfont
\renewcommand{\arraystretch}{1.15}
\setlength{\tabcolsep}{6pt}
\begin{longtable}{|>{\raggedright\arraybackslash}p{0.09\textwidth}|>{\raggedright\arraybackslash}p{0.84\textwidth}|}
\hline
\textbf{Field} & \textbf{Value} \\ \hline
\endhead
Use Case ID & UC-06 \\ \hline
Name & Share Learning Resource \\ \hline
Primary Actor & University Student \\ \hline
Supporting Actors & None \\ \hline
Feature Group & Learning Resource Library \\ \hline
Priority & High \\ \hline
Related User Stories & US-27 \\ \hline
Related Requirements & US-27, US-25, US-28, US-29; FR-27.1 to FR-27.5, FR-25.1 to FR-25.5, FR-28.1 to FR-28.3, FR-29.1 to FR-29.4; NFR-SEC-5, NFR-PRIV-5, NFR-SCAL-4 \\ \hline
Source Documents & \texttt{docs/requirements/proposals/use\_cases\_model/UC-06-share-learning-resource.md}; \texttt{docs/requirements/proposals/use\_cases\_description.md} (UC-06); \texttt{docs/requirements/proposals/use\_cases\_model.md}; \texttt{docs/requirements/proposals/use\_case\_diagrams.md}; \texttt{docs/requirements/PA3-Group05/use\_case.md} (section 6 UC-06); \texttt{docs/requirements/proposals/overview.md}; \texttt{docs/requirements/proposals/functional\_requirement.md}; \texttt{docs/requirements/proposals/non\_functional\_requirement.md}; \texttt{docs/requirements/proposals/userstory.md}; \texttt{docs/requirements/proposals/feature\_flow.md}; \texttt{docs/requirements/proposals/screen\_flow.md}; \texttt{docs/requirements/proposals/API.md}; \texttt{docs/requirements/proposals/response\_code.md}; \texttt{docs/requirements/proposals/library.md}; \texttt{docs/requirements/proposals/create\_post.md}; \texttt{docs/requirements/proposals/group\_detail.md}; \texttt{docs/requirements/proposals/friend.md}; \texttt{docs/requirements/proposals/friend\_profile.md}; \texttt{docs/requirements/proposals/settings.md} \\ \hline
\end{longtable}
\endgroup
\par\medskip

\hypertarget{sec:part-a-use-case-model}{}
\addcontentsline{toc}{subsection}{Part A --- Use Case Model}
\subsection*{Part A --- Use Case Model}
\par\medskip

\hypertarget{sec:a1-focused-use-case-diagram}{}
\addcontentsline{toc}{subsubsection}{A1. Focused Use Case Diagram}
\subsubsection*{A1. Focused Use Case Diagram}
\par\medskip

\begin{center}
\includegraphics[width=0.85\textwidth,height=0.75\textheight,keepaspectratio]{"figures/UC-06-share-learning-resource_1.png"}
\end{center}
\par\medskip

\textbf{Diagram notation.} Person shape = human actor; blue = selected base use case; purple = included; green = extending; solid line = association; dashed arrow = include or extend.
\par\medskip

\hypertarget{sec:a2-actor-register}{}
\addcontentsline{toc}{subsubsection}{A2. Actor Register}
\subsubsection*{A2. Actor Register}
\par\medskip

\begingroup
\small
\renewcommand{\baselinestretch}{1.0}\selectfont
\renewcommand{\arraystretch}{1.15}
\setlength{\tabcolsep}{6pt}
\begin{longtable}{|>{\raggedright\arraybackslash}p{0.36\textwidth}|>{\raggedright\arraybackslash}p{0.10\textwidth}|>{\raggedright\arraybackslash}p{0.46\textwidth}|}
\hline
\textbf{Actor} & \textbf{Type} & \textbf{Description} \\ \hline
\endhead
University Student & Primary & Authenticated student who publishes a learning resource to the Learning Resource Library for discovery and reuse. \\ \hline
New University Student; Registered University Student; Post Author; Room Owner or Manager & Specialized human actors & Specialized student actors that inherit the University Student association to UC-06 in the shared actor model. \\ \hline
\end{longtable}
\endgroup
\par\medskip

\hypertarget{sec:a3-relationships}{}
\addcontentsline{toc}{subsubsection}{A3. Relationships}
\subsubsection*{A3. Relationships}
\par\medskip

\hypertarget{sec:a3-1-associations}{}
\paragraph*{A3.1 Associations}
\par\medskip

\begingroup
\small
\renewcommand{\baselinestretch}{1.0}\selectfont
\renewcommand{\arraystretch}{1.15}
\setlength{\tabcolsep}{6pt}
\begin{longtable}{|>{\raggedright\arraybackslash}p{0.15\textwidth}|>{\raggedright\arraybackslash}p{0.23\textwidth}|>{\raggedright\arraybackslash}p{0.10\textwidth}|>{\raggedright\arraybackslash}p{0.44\textwidth}|}
\hline
\textbf{Actor} & \textbf{Use Case} & \textbf{Type} & \textbf{Description} \\ \hline
\endhead
University Student & UC-06 Share Learning Resource & Primary & Publishes a learning resource to the resource library. \\ \hline
\end{longtable}
\endgroup
\par\medskip

\hypertarget{sec:a3-2-generalization-and-specialization}{}
\paragraph*{A3.2 Generalization and Specialization}
\par\medskip

\begingroup
\small
\renewcommand{\baselinestretch}{1.0}\selectfont
\renewcommand{\arraystretch}{1.15}
\setlength{\tabcolsep}{6pt}
\begin{longtable}{|>{\raggedright\arraybackslash}p{0.42\textwidth}|>{\raggedright\arraybackslash}p{0.10\textwidth}|>{\raggedright\arraybackslash}p{0.10\textwidth}|>{\raggedright\arraybackslash}p{0.29\textwidth}|}
\hline
\textbf{Child} & \textbf{Parent} & \textbf{Type} & \textbf{Inherited Behavior or Association} \\ \hline
\endhead
New University Student; Registered University Student; Post Author; Room Owner or Manager & University Student & Actor generalization & Each specialized student actor inherits the UC-06 association. \\ \hline
\end{longtable}
\endgroup
\par\medskip

\hypertarget{sec:a3-3-include}{}
\paragraph*{A3.3 Include}
\par\medskip

\begingroup
\small
\renewcommand{\baselinestretch}{1.0}\selectfont
\renewcommand{\arraystretch}{1.15}
\setlength{\tabcolsep}{6pt}
\begin{longtable}{|>{\raggedright\arraybackslash}p{0.19\textwidth}|>{\raggedright\arraybackslash}p{0.21\textwidth}|>{\raggedright\arraybackslash}p{0.52\textwidth}|}
\hline
\textbf{Base Use Case} & \textbf{Included Use Case} & \textbf{Meaning} \\ \hline
\endhead
UC-06 Share Learning Resource & I-07 Validate and Store Resource & Validate the file and metadata, store the resource, and index it for search. \\ \hline
\end{longtable}
\endgroup
\par\medskip

\hypertarget{sec:a3-4-extend}{}
\paragraph*{A3.4 Extend}
\par\medskip

\begingroup
\small
\renewcommand{\baselinestretch}{1.0}\selectfont
\renewcommand{\arraystretch}{1.15}
\setlength{\tabcolsep}{6pt}
\begin{longtable}{|>{\raggedright\arraybackslash}p{0.15\textwidth}|>{\raggedright\arraybackslash}p{0.15\textwidth}|>{\raggedright\arraybackslash}p{0.25\textwidth}|>{\raggedright\arraybackslash}p{0.35\textwidth}|}
\hline
\textbf{Extending Use Case} & \textbf{Base Use Case} & \textbf{Condition} & \textbf{Added Behavior} \\ \hline
\endhead
E-07 Restrict Resource Access & UC-06 Share Learning Resource & The owner chooses a restricted privacy setting. & Limit resource visibility and access according to the privacy rule. \\ \hline
\end{longtable}
\endgroup
\par\medskip

\hypertarget{sec:a4-model-notes}{}
\addcontentsline{toc}{subsubsection}{A4. Model Notes}
\subsubsection*{A4. Model Notes}
\par\medskip

\begin{enumerate}
\item Every published resource includes I-07 because the resource must be validated, persisted, and indexed before it is available.
\item Restricting resource access is an optional extension of the base sharing behavior.
\item UC-06 belongs to the Learning Resource Library feature group and remains part of the primary student workflow.
\item The focused model keeps supporting actors out of scope; file storage and file-safety checks are system/service preconditions in the specification, not modeled supporting actors.
\end{enumerate}
\par\medskip

\hypertarget{sec:part-b-use-case-specification}{}
\addcontentsline{toc}{subsection}{Part B --- Use Case Specification}
\subsection*{Part B --- Use Case Specification}
\par\medskip

\hypertarget{sec:b1-brief-description}{}
\addcontentsline{toc}{subsubsection}{B1. Brief Description}
\subsubsection*{B1. Brief Description}
\par\medskip

A student publishes a learning resource so that other students can discover and access useful academic material. The resource includes metadata, optional tags, and an uploaded file, and becomes searchable after publication.
\par\medskip

\hypertarget{sec:b2-preconditions}{}
\addcontentsline{toc}{subsubsection}{B2. Preconditions}
\subsubsection*{B2. Preconditions}
\par\medskip

\begin{enumerate}
\item The student is authenticated with an active session.
\item The student is permitted to publish resources to the library.
\item The file storage service and configured file-validation checks are available.
\end{enumerate}
\par\medskip

\hypertarget{sec:b3-postconditions}{}
\addcontentsline{toc}{subsubsection}{B3. Postconditions}
\subsubsection*{B3. Postconditions}
\par\medskip

\begin{enumerate}
\item \textbf{Success:} A validated resource is stored with its owner, metadata, file, subject tags, and course tags.
\item \textbf{Success:} The resource is searchable in the library after publication.
\item \textbf{Success:} Other users can view and access the resource according to its privacy settings.
\item \textbf{Failure:} No resource is published, and the student receives a clear validation or system error.
\end{enumerate}
\par\medskip

\hypertarget{sec:b4-main-success-scenario-basic-flow}{}
\addcontentsline{toc}{subsubsection}{B4. Main Success Scenario (Basic Flow)}
\subsubsection*{B4. Main Success Scenario (Basic Flow)}
\par\medskip

\begin{enumerate}
\item The student opens the Learning Resource Library and selects \textbf{Share Resource}.
\item The system displays the resource form for title, description, subject tags, course tags, and file upload.
\item The student enters the resource details and selects a file.
\item The frontend validates the required fields and checks the file type and size against the configured limits.
\item The student submits the resource.
\item The backend verifies authentication and publication permission, validates the metadata and file, and performs the configured file-safety checks.
\item The backend stores the file and creates the resource record associated with the student, tags, and academic context.
\item The system publishes the resource and makes it searchable in the library.
\item The system displays a success confirmation and the resource in the student's storage view.
\item The resource becomes available to other users according to its privacy settings.
\end{enumerate}
\par\medskip

\hypertarget{sec:b5-alternative-flows}{}
\addcontentsline{toc}{subsubsection}{B5. Alternative Flows}
\subsubsection*{B5. Alternative Flows}
\par\medskip

Unless otherwise stated, an alternative flow resumes at the next step of the basic flow.
\par\medskip

\hypertarget{sec:a1-add-or-change-tags}{}
\paragraph*{A1. Add or change tags}
\par\medskip

At basic flow step 3, the student adds or removes subject or course tags. The system updates the searchable metadata before publication.
\par\medskip

\hypertarget{sec:a2-replace-the-file-before-submission}{}
\paragraph*{A2. Replace the file before submission}
\par\medskip

The student selects a different file. The system revalidates file type and size and replaces the pending upload.
\par\medskip

\hypertarget{sec:a3-privacy-restricted-resource}{}
\paragraph*{A3. Privacy-restricted resource}
\par\medskip

At basic flow step 8, the resource is published with restricted access. It becomes searchable according to the privacy rule, but only authorized users can access the file.
\par\medskip

\hypertarget{sec:a4-cancel-publication}{}
\paragraph*{A4. Cancel publication}
\par\medskip

The student leaves the form before submitting. No published resource is created.
\par\medskip

\hypertarget{sec:a5-view-own-published-resources}{}
\paragraph*{A5. View own published resources}
\par\medskip

After publication, the student opens the storage view to see the resource status, publication date, and access statistics.
\par\medskip

\hypertarget{sec:b6-exception-flows}{}
\addcontentsline{toc}{subsubsection}{B6. Exception Flows}
\subsubsection*{B6. Exception Flows}
\par\medskip

\hypertarget{sec:e1-missing-or-invalid-metadata}{}
\paragraph*{E1. Missing or invalid metadata}
\par\medskip

The system highlights the invalid fields, keeps the form open, and does not publish.
\par\medskip

\hypertarget{sec:e2-unsupported-file-type-or-oversized-file}{}
\paragraph*{E2. Unsupported file type or oversized file}
\par\medskip

The system rejects the file, identifies the violated constraint, and allows the student to choose another file.
\par\medskip

\hypertarget{sec:e3-file-safety-check-fails}{}
\paragraph*{E3. File-safety check fails}
\par\medskip

The system blocks publication and informs the student that the file cannot be accepted.
\par\medskip

\hypertarget{sec:e4-storage-or-database-failure}{}
\paragraph*{E4. Storage or database failure}
\par\medskip

No published resource is created. The system displays a retry option and avoids leaving an orphaned or partially published resource.
\par\medskip

\hypertarget{sec:e5-publication-permission-lost}{}
\paragraph*{E5. Publication permission lost}
\par\medskip

The system rejects the publication and explains that the student is no longer authorized.
\par\medskip

\hypertarget{sec:e6-resource-becomes-unavailable-after-publication}{}
\paragraph*{E6. Resource becomes unavailable after publication}
\par\medskip

The system prevents further access and removes the resource from search and library listings when required.
\par\medskip

\hypertarget{sec:b7-business-rules}{}
\addcontentsline{toc}{subsubsection}{B7. Business Rules}
\subsubsection*{B7. Business Rules}
\par\medskip

\begin{enumerate}
\item A resource must include a title and the metadata required for publication.
\item A resource can include a description, subject tags, and course tags.
\item Uploaded files must pass file type and size validation before acceptance.
\item A published resource must be searchable in the library.
\item Access to a published resource must follow the owner's privacy settings.
\item A deleted resource must be removed from search and library listings and must no longer be accessible to other users.
\item Publication is restricted to authenticated and authorized students.
\end{enumerate}
\par\medskip

\hypertarget{sec:b8-special-requirements}{}
\addcontentsline{toc}{subsubsection}{B8. Special Requirements}
\subsubsection*{B8. Special Requirements}
\par\medskip

\begingroup
\small
\renewcommand{\baselinestretch}{1.0}\selectfont
\renewcommand{\arraystretch}{1.15}
\setlength{\tabcolsep}{6pt}
\begin{longtable}{|>{\raggedright\arraybackslash}p{0.10\textwidth}|>{\raggedright\arraybackslash}p{0.83\textwidth}|}
\hline
\textbf{ID} & \textbf{Requirement} \\ \hline
\endhead
NFR-SEC-5 & Uploaded files must be validated for file type and file size before acceptance. \\ \hline
NFR-PRIV-5 & Users must be able to control access to documents and learning resources they share. \\ \hline
NFR-SCAL-4 & The architecture must support at least 100,000 uploaded documents and learning resources. \\ \hline
Design constraint & Configured file-safety checks must complete before publication, and failed storage or database operations must not leave orphaned or partially published resources. \\ \hline
\end{longtable}
\endgroup
\par\medskip

\hypertarget{sec:b9-extension-points}{}
\addcontentsline{toc}{subsubsection}{B9. Extension Points}
\subsubsection*{B9. Extension Points}
\par\medskip

\begingroup
\small
\renewcommand{\baselinestretch}{1.0}\selectfont
\renewcommand{\arraystretch}{1.15}
\setlength{\tabcolsep}{6pt}
\begin{longtable}{|>{\raggedright\arraybackslash}p{0.36\textwidth}|>{\raggedright\arraybackslash}p{0.57\textwidth}|}
\hline
\textbf{Extension Point} & \textbf{Location in Flow of Events} \\ \hline
\endhead
Tag changes & Basic flow step 3; alternative A1 \\ \hline
File replacement & Basic flow step 3; alternative A2 \\ \hline
Privacy restriction & Basic flow step 8; alternative A3 \\ \hline
Metadata validation failure & Basic flow step 4; exception E1 \\ \hline
File validation failure & Basic flow steps 4 and 6; exceptions E2-E3 \\ \hline
Cancel publication & Before basic flow step 5; alternative A4 \\ \hline
Storage-view review & Basic flow step 9; alternative A5 \\ \hline
\end{longtable}
\endgroup
\par\medskip

\hypertarget{sec:b10-interface-and-data-notes}{}
\addcontentsline{toc}{subsubsection}{B10. Interface and Data Notes}
\subsubsection*{B10. Interface and Data Notes}
\par\medskip

\begin{itemize}
\item \textbf{Screen flow and entry points:} \texttt{screen\_flow.md} places the library in the main flow as \texttt{Home -$>$ Library Group}, with \texttt{Library -$>$ Exam} and \texttt{Library -$>$ Book}. The library page recommendation uses \texttt{LibraryPage} / \texttt{LearningResourceLibraryPage}, tabbed views (\textbf{All}, \textbf{Books}, \textbf{Exams}, \textbf{Recommended}, \textbf{Saved}, \textbf{My Storage}), and a primary \textbf{Share Resource} action.
\item \textbf{Share form fields:} \texttt{ShareResourceDialog} / \texttt{ShareResourceForm} captures \texttt{ResourceTitleInput} (required title), \texttt{ResourceDescriptionInput}, \texttt{ResourceSubjectTagInput}, \texttt{ResourceCourseTagInput}, \texttt{ResourceAcademicContextSelector} (University, Major, Course), \texttt{ResourceFileDropzone}, \texttt{ResourceFilePreviewPanel}, and \texttt{ResourcePrivacySelect}.
\item \textbf{Privacy levels and meanings:} \texttt{Public} = any authenticated Studitopia user can view and access the resource; \texttt{University} = only users from the owner's university can access the file; \texttt{Course} = only members of the selected course can access the file; \texttt{Private} = only the owner can access the file. Access must follow the owner's privacy setting (FR-27.5, NFR-PRIV-5, E-07 Restrict Resource Access).
\item \textbf{Validation and file-safety rules:} Frontend validates required fields and checks file type/size before submission (FR-27.3). Backend verifies authentication/publication permission, validates metadata and file, and performs configured file-safety checks (NFR-SEC-5). Unsupported file type, oversized file, missing/invalid metadata, or failed file-safety checks block publication and keep the form actionable.
\item \textbf{Search/indexing:} A published resource becomes searchable in the library after publication (FR-27.4, AC-27.1). Search/filter support includes title, author, subject tags, description, format, subject field, and academic context (FR-25.1 to FR-25.5); ratings/review counts may be displayed when available (FR-25.4).
\item \textbf{Exact UI messages and labels:} \textbf{Share Resource} (primary action), \textbf{Share a learning resource} (dialog heading), \textbf{Publish resource} (submit), \textbf{Cancel} (cancel), \textbf{Retry} (retry after failure), “Resource published. It is now searchable in the library.” (\texttt{PublishSuccessToast}), “No resources found” (\texttt{NoResultsState}), “Unable to load resources.” with \textbf{Retry} (\texttt{LibraryErrorBanner}), “Unsupported file type” / “File exceeds maximum size.” (\texttt{FileValidationErrorText}), “You do not have permission to access this resource.” (\texttt{AccessDeniedState}), “Resource could not be saved. Please retry.” (\texttt{StorageFailureAlert}), and “You have not published any resources yet.” (\texttt{MyStorageEmptyState}).
\item \textbf{API/material-post path and codes:} \texttt{API.md} defines no standalone library upload endpoint. The documented material-sharing path is \texttt{POST /forum/posts} with \texttt{post\_type: 'material'}, \texttt{title}, \texttt{content}, optional \texttt{tags}, and \texttt{files}; response \texttt{201} returns \texttt{post\_id} and \texttt{created\_at}. Related forum error codes include \texttt{UNKNOWN\_ERROR: 45000}, \texttt{INVALID\_POST\_TYPE: 48003}, \texttt{INVALID\_POST\_CONTENT: 48004}, \texttt{FORBIDDEN: 48010}, and \texttt{POST\_NOT\_FOUND: 48001} where applicable to forum/material flows. Library-specific upload, privacy, file-safety, storage, and indexing endpoints/codes are not specified in \texttt{API.md}/\texttt{response\_code.md}.
\item \textbf{Group/friend material-sharing rules:} Friend and group flows can share materials/documents in messages (\texttt{MaterialShareButton}, \texttt{MessageAttachmentButton}, \texttt{GroupMaterialShareButton}, \texttt{SharedDocumentPanel}). Group detail states published/shared resources follow privacy settings and access is allowed only when permitted; material may also appear in the Learning Resource Library where applicable (\texttt{FR-15.6}, \texttt{FR-27.4}, \texttt{FR-27.5}, \texttt{NFR-PRIV-5}, \texttt{NFR-PRIV-7}). Library publication remains the UC-06 resource-sharing path; group/friend sharing must respect the same privacy and access constraints.
\end{itemize}
\par\medskip

\hypertarget{sec:b11-source-alignment-notes}{}
\addcontentsline{toc}{subsubsection}{B11. Source Alignment Notes}
\subsubsection*{B11. Source Alignment Notes}
\par\medskip

\begin{itemize}
\item \textbf{UC-06 source alignment:} Main flow, alternatives A1-A5, exceptions E1-E6, business rules, and extension points are aligned to \texttt{use\_cases\_description.md} UC-06 and \texttt{PA3-Group05/use\_case.md} section 6.
\item \textbf{Diagram fidelity:} The A1 Mermaid diagram is copied verbatim from the focused UC-06 model file; relationship tables follow the focused model structure and the shared model register.
\item \textbf{API gap:} \texttt{API.md} has no standalone library upload endpoint (only the material post path \texttt{POST /forum/posts} with \texttt{post\_type: 'material'}). This document does not invent library upload routes, request schemas, or error codes.
\item \textbf{Privacy and storage details:} Privacy labels/meanings, storage statistics, and some UI component names are taken from \texttt{library.md} as UI recommendations consistent with FR-27/FR-28/FR-29 and NFR-PRIV-5; exact backend enums and file constraints are not specified in the API proposal.
\item \textbf{Cross-feature consistency:} Material posts may appear in the library when applicable (\texttt{FR-15.6}), and friend/group material sharing must follow privacy/access rules, but those flows are not substitutes for UC-06 library publication.
\end{itemize}
\par\medskip

\par\medskip
\hypertarget{sec:ui-prototype-specimens}{}
\addcontentsline{toc}{subsection}{UI Prototype Specimens}
\subsection*{UI Prototype Specimens}
\par\medskip
\begin{figure}[H]
\centering
\includegraphics[width=0.75\textwidth,height=0.42\textheight,keepaspectratio]{"../../prototype/create_post_material.png"}
\par\vspace{2pt}
{\footnotesize \textbf{UI Prototype:} Create post - Material}
\end{figure}
\par\smallskip

\pagebreak
\hypertarget{sec:uc-22-find-learning-resources}{}
\addcontentsline{toc}{section}{UC-22 --- Find Learning Resources}
\section*{UC-22 --- Find Learning Resources}
\par\medskip

\textbf{Version:} \texttt{1.0}
\par\medskip

\textbf{Performed by:} [Nguyen Le Quoc Huy] | \textbf{Edited by:} [Nguyen Le Quoc Huy]
\par\medskip

\begingroup
\small
\renewcommand{\baselinestretch}{1.0}\selectfont
\renewcommand{\arraystretch}{1.15}
\setlength{\tabcolsep}{6pt}
\begin{longtable}{|>{\raggedright\arraybackslash}p{0.09\textwidth}|>{\raggedright\arraybackslash}p{0.84\textwidth}|}
\hline
\textbf{Field} & \textbf{Value} \\ \hline
\endhead
Use Case ID & UC-22 \\ \hline
Name & Find Learning Resources \\ \hline
Primary Actor & University Student \\ \hline
Supporting Actors & None \\ \hline
Feature Group & Learning Resource Library \\ \hline
Priority & Not explicitly stated \\ \hline
Related User Stories & US-25 \\ \hline
Related Requirements & FR-25.1 to FR-25.5; AC-25.1 to AC-25.4; NFR-PERF-1; NFR-SCAL-4; NFR-SEC-1; NFR-SEC-6; NFR-PRIV-5; NFR-PRIV-7; NFR-USE-2; NFR-USE-4 \\ \hline
Source Documents & \texttt{docs/requirements/proposals/userstory.md}; \texttt{docs/requirements/proposals/functional\_requirement.md}; \texttt{docs/requirements/proposals/library.md}; \texttt{docs/requirements/proposals/create\_post.md}; \texttt{docs/requirements/proposals/API.md}; \texttt{docs/requirements/proposals/response\_code.md}; \texttt{docs/requirements/proposals/non\_functional\_requirement.md}; \texttt{docs/requirements/proposals/feature\_flow.md}; \texttt{docs/requirements/proposals/screen\_flow.md} \\ \hline
\end{longtable}
\endgroup
\par\medskip

\hypertarget{sec:part-a-use-case-model}{}
\addcontentsline{toc}{subsection}{Part A --- Use Case Model}
\subsection*{Part A --- Use Case Model}
\par\medskip

\hypertarget{sec:a1-focused-use-case-diagram}{}
\addcontentsline{toc}{subsubsection}{A1. Focused Use Case Diagram}
\subsubsection*{A1. Focused Use Case Diagram}
\par\medskip

\begin{center}
\includegraphics[width=0.85\textwidth,height=0.75\textheight,keepaspectratio]{"figures/UC-22-find-learning-resources_1.png"}
\end{center}
\par\medskip

\textbf{Diagram notation.} Person shape = human actor; cloud = supporting system; blue = selected use case; solid line = association; labelled solid arrow = generalization; dashed arrow = include or extend.
\par\medskip

\hypertarget{sec:a2-actor-register}{}
\addcontentsline{toc}{subsubsection}{A2. Actor Register}
\subsubsection*{A2. Actor Register}
\par\medskip

\begingroup
\small
\renewcommand{\baselinestretch}{1.0}\selectfont
\renewcommand{\arraystretch}{1.15}
\setlength{\tabcolsep}{6pt}
\begin{longtable}{|>{\raggedright\arraybackslash}p{0.13\textwidth}|>{\raggedright\arraybackslash}p{0.09\textwidth}|>{\raggedright\arraybackslash}p{0.69\textwidth}|}
\hline
\textbf{Actor} & \textbf{Type} & \textbf{Description} \\ \hline
\endhead
University Student & Primary & Authenticated student who searches books and learning resources related to the student's courses. \\ \hline
\end{longtable}
\endgroup
\par\medskip

\hypertarget{sec:a3-relationships}{}
\addcontentsline{toc}{subsubsection}{A3. Relationships}
\subsubsection*{A3. Relationships}
\par\medskip

\hypertarget{sec:a3-1-associations}{}
\paragraph*{A3.1 Associations}
\par\medskip

\begingroup
\small
\renewcommand{\baselinestretch}{1.0}\selectfont
\renewcommand{\arraystretch}{1.15}
\setlength{\tabcolsep}{6pt}
\begin{longtable}{|>{\raggedright\arraybackslash}p{0.13\textwidth}|>{\raggedright\arraybackslash}p{0.21\textwidth}|>{\raggedright\arraybackslash}p{0.09\textwidth}|>{\raggedright\arraybackslash}p{0.47\textwidth}|}
\hline
\textbf{Actor} & \textbf{Use Case} & \textbf{Type} & \textbf{Description} \\ \hline
\endhead
University Student & UC-22 Find Learning Resources & Primary & Searches, filters, and reviews learning resources in the library. \\ \hline
\end{longtable}
\endgroup
\par\medskip

\hypertarget{sec:a3-2-generalization-and-specialization}{}
\paragraph*{A3.2 Generalization and Specialization}
\par\medskip

\begingroup
\small
\renewcommand{\baselinestretch}{1.0}\selectfont
\renewcommand{\arraystretch}{1.15}
\setlength{\tabcolsep}{6pt}
\begin{longtable}{|>{\raggedright\arraybackslash}p{0.09\textwidth}|>{\raggedright\arraybackslash}p{0.09\textwidth}|>{\raggedright\arraybackslash}p{0.26\textwidth}|>{\raggedright\arraybackslash}p{0.47\textwidth}|}
\hline
\textbf{Child} & \textbf{Parent} & \textbf{Type} & \textbf{Inherited Behavior or Association} \\ \hline
\endhead
None & None & No generalization or specialization & The proposals name only the University Student actor for US-25. \\ \hline
\end{longtable}
\endgroup
\par\medskip

\hypertarget{sec:a3-3-include}{}
\paragraph*{A3.3 Include}
\par\medskip

\begingroup
\small
\renewcommand{\baselinestretch}{1.0}\selectfont
\renewcommand{\arraystretch}{1.15}
\setlength{\tabcolsep}{6pt}
\begin{longtable}{|>{\raggedright\arraybackslash}p{0.23\textwidth}|>{\raggedright\arraybackslash}p{0.30\textwidth}|>{\raggedright\arraybackslash}p{0.40\textwidth}|}
\hline
\textbf{Base Use Case} & \textbf{Included Use Case} & \textbf{Meaning} \\ \hline
\endhead
None & None & No include relationship & The proposals do not name an included use case for US-25. \\ \hline
\end{longtable}
\endgroup
\par\medskip

\hypertarget{sec:a3-4-extend}{}
\paragraph*{A3.4 Extend}
\par\medskip

\begingroup
\small
\renewcommand{\baselinestretch}{1.0}\selectfont
\renewcommand{\arraystretch}{1.15}
\setlength{\tabcolsep}{6pt}
\begin{longtable}{|>{\raggedright\arraybackslash}p{0.15\textwidth}|>{\raggedright\arraybackslash}p{0.11\textwidth}|>{\raggedright\arraybackslash}p{0.18\textwidth}|>{\raggedright\arraybackslash}p{0.47\textwidth}|}
\hline
\textbf{Extending Use Case} & \textbf{Base Use Case} & \textbf{Condition} & \textbf{Added Behavior} \\ \hline
\endhead
None & None & No extend relationship & The proposals do not name an extending use case for US-25. \\ \hline
\end{longtable}
\endgroup
\par\medskip

\hypertarget{sec:a4-model-notes}{}
\addcontentsline{toc}{subsubsection}{A4. Model Notes}
\subsubsection*{A4. Model Notes}
\par\medskip

\begin{enumerate}
\item The scope is keyword search, field-specific search, filtering, sorting, ratings display, and academic relevance in the Learning Resource Library.
\item Publishing, deleting, and managing storage are outside this use case; they are separate learning-resource stories.
\item No distinct included, extending, or specialized use case is named for US-25, so no I-xx, E-xx, or G-xx relationship ID is introduced.
\item The focused model shows only the primary human actor because the proposals do not name a supporting system actor for this story.
\item The exact API endpoint, route, result limit, and priority are not specified in proposals.
\end{enumerate}
\par\medskip

\hypertarget{sec:part-b-use-case-specification}{}
\addcontentsline{toc}{subsection}{Part B --- Use Case Specification}
\subsection*{Part B --- Use Case Specification}
\par\medskip

\hypertarget{sec:b1-brief-description}{}
\addcontentsline{toc}{subsubsection}{B1. Brief Description}
\subsubsection*{B1. Brief Description}
\par\medskip

A university student searches for books and learning resources so that the student can find materials related to courses. The search covers title, author, subject tags, and description, with filtering by format and subject field; results may also show aggregate ratings and review counts when available.
\par\medskip

The library reinforces the forum by preserving and reusing academic material. Material posts created in the forum may also appear in the library when applicable (\texttt{FR-15.6}), and the library displays resources related to the user's university, major, and courses (\texttt{AC-25.4}, \texttt{FR-25.5}).
\par\medskip

\hypertarget{sec:b2-preconditions}{}
\addcontentsline{toc}{subsubsection}{B2. Preconditions}
\subsubsection*{B2. Preconditions}
\par\medskip

\begin{enumerate}
\item The student is authenticated before accessing protected user resources (\texttt{NFR-SEC-1}).
\item The Learning Resource Library is reachable in the main flow as \texttt{Home -$>$ Library Group}, with \texttt{Library Group -$>$ Exam} and \texttt{Library Group -$>$ Book} from \texttt{screen\_flow.md}.
\item Published resources are available for search after publication (\texttt{FR-27.4}, \texttt{AC-27.1}).
\item University, major, course, subject tags, and privacy data are available for relevance and access decisions (\texttt{FR-25.5}, \texttt{NFR-PRIV-7}).
\end{enumerate}
\par\medskip

\hypertarget{sec:b3-postconditions}{}
\addcontentsline{toc}{subsubsection}{B3. Postconditions}
\subsubsection*{B3. Postconditions}
\par\medskip

\begin{enumerate}
\item \textbf{Success:} Matching resources are displayed with the requested search terms, filters, and sort order.
\item \textbf{Success:} The result list shows ratings and review counts when available and shows resources related to the student's university, major, and courses.
\item \textbf{Success:} The student can open a result for detail or save it for later where the library UI supports that action.
\item \textbf{Failure:} No misleading result is shown, and the student receives the specified empty-state or load-error feedback.
\end{enumerate}
\par\medskip

\hypertarget{sec:b4-main-success-scenario-basic-flow}{}
\addcontentsline{toc}{subsubsection}{B4. Main Success Scenario (Basic Flow)}
\subsubsection*{B4. Main Success Scenario (Basic Flow)}
\par\medskip

\begin{enumerate}
\item The student opens the library from Home, and the system displays the Learning Resource Library page.
\item The student selects the search area represented by \texttt{ResourceSearchPanel} with heading \textbf{Find a resource}.
\item The student enters keywords through \texttt{ResourceKeywordSearchInput}, whose exact placeholder is “Search title, author, tag, or description.”
\item The student selects \textbf{Search} to submit the keyword query.
\item The system searches by title, author, subject tags, and description (\texttt{AC-25.1}, \texttt{FR-25.2}).
\item The student selects filters for format and subject field (\texttt{AC-25.2}, \texttt{FR-25.3}).
\item The student selects a sort option from \textbf{Relevance}, \textbf{Top rated}, \textbf{Most accessed}, and \textbf{Newest}.
\item The system displays matching \texttt{ResourceCard} items in \texttt{ResourceResultGrid} with a \texttt{ResultCountText}.
\item The system displays aggregate ratings and review counts when available (\texttt{AC-25.3}, \texttt{FR-25.4}).
\item The system displays resources related to the student's university, major, and courses (\texttt{AC-25.4}, \texttt{FR-25.5}).
\item The student selects \textbf{View details} to inspect a resource, or uses \textbf{Save} where the library UI supports saving.
\end{enumerate}
\par\medskip

\hypertarget{sec:b5-alternative-flows}{}
\addcontentsline{toc}{subsubsection}{B5. Alternative Flows}
\subsubsection*{B5. Alternative Flows}
\par\medskip

\hypertarget{sec:a1-advanced-field-specific-search}{}
\paragraph*{A1. Advanced field-specific search}
\par\medskip

The student opens \textbf{Advanced filters} and narrows results with \textbf{Title}, \textbf{Author}, subject-tag chips such as \texttt{Calculus}, \texttt{Database}, and \texttt{Physics}, or \textbf{Description contains}. The system applies these as field-specific refinements of the keyword query (\texttt{FR-25.2}).
\par\medskip

\hypertarget{sec:a2-course-and-academic-scope-filtering}{}
\paragraph*{A2. Course and academic-scope filtering}
\par\medskip

The student selects course tags such as \texttt{CS101} and \texttt{MATH202}, or a University, Major, Course scope through \texttt{ResourceAcademicScopeFilter}. The system limits results to the selected academic context (\texttt{FR-25.5}, \texttt{AC-25.4}).
\par\medskip

\hypertarget{sec:a3-ratings-unavailable}{}
\paragraph*{A3. Ratings unavailable}
\par\medskip

If a resource has no rating data, the system omits \texttt{ResourceRatingDisplay} or shows “Not yet rated.” instead of inventing a rating (\texttt{FR-25.4}, \texttt{library.md} business rules).
\par\medskip

\hypertarget{sec:a4-clear-filters}{}
\paragraph*{A4. Clear filters}
\par\medskip

The student removes individual chips from \texttt{ActiveFilterChipList} or selects \textbf{Clear all filters}. The system resets the search and filters and recalculates \texttt{ResultCountText}.
\par\medskip

\hypertarget{sec:b6-exception-flows}{}
\addcontentsline{toc}{subsubsection}{B6. Exception Flows}
\subsubsection*{B6. Exception Flows}
\par\medskip

\hypertarget{sec:e1-no-results}{}
\paragraph*{E1. No results}
\par\medskip

No resource matches the criteria. The system displays the exact empty-state message “No resources found” and offers actions to broaden filters or share a resource.
\par\medskip

\hypertarget{sec:e2-resource-load-failure}{}
\paragraph*{E2. Resource load failure}
\par\medskip

The library cannot load results. The system shows “Unable to load resources.” with \textbf{Retry} through \texttt{LibraryErrorBanner}.
\par\medskip

\hypertarget{sec:e3-privacy-blocked-access}{}
\paragraph*{E3. Privacy-blocked access}
\par\medskip

The student attempts to access a resource whose privacy setting does not permit access. The system explains the restriction with “You do not have permission to access this resource.” and does not expose the file.
\par\medskip

\hypertarget{sec:b7-business-rules}{}
\addcontentsline{toc}{subsubsection}{B7. Business Rules}
\subsubsection*{B7. Business Rules}
\par\medskip

\begin{enumerate}
\item Keywords search title, author, subject tags, and description (\texttt{AC-25.1}, \texttt{FR-25.2}).
\item Results can be filtered by format and subject field (\texttt{AC-25.2}, \texttt{FR-25.3}).
\item Search results display aggregate ratings and review counts when available (\texttt{AC-25.3}, \texttt{FR-25.4}).
\item The library displays resources related to the user's university, major, and courses (\texttt{AC-25.4}, \texttt{FR-25.5}).
\item Published resources become searchable in the library after publication (\texttt{FR-27.4}, \texttt{AC-27.1}).
\item Access and search visibility must follow owner privacy settings (\texttt{FR-27.5}, \texttt{NFR-PRIV-5}, \texttt{NFR-PRIV-7}).
\end{enumerate}
\par\medskip

\hypertarget{sec:b8-special-requirements}{}
\addcontentsline{toc}{subsubsection}{B8. Special Requirements}
\subsubsection*{B8. Special Requirements}
\par\medskip

\begingroup
\small
\renewcommand{\baselinestretch}{1.0}\selectfont
\renewcommand{\arraystretch}{1.15}
\setlength{\tabcolsep}{6pt}
\begin{longtable}{|>{\raggedright\arraybackslash}p{0.10\textwidth}|>{\raggedright\arraybackslash}p{0.83\textwidth}|}
\hline
\textbf{ID} & \textbf{Requirement} \\ \hline
\endhead
NFR-PERF-1 & The system shall provide responsive user interactions with a response time of ≤ 2 seconds for standard user actions. \\ \hline
NFR-SCAL-4 & The system architecture shall support at least 100,000 uploaded documents and learning resources. \\ \hline
NFR-SEC-1 & The system shall require authentication before accessing protected user resources. \\ \hline
NFR-SEC-6 & The system shall prevent unauthorized access through server-side authorization checks. \\ \hline
NFR-PRIV-5 & Users shall be able to control access to documents and learning resources they share. \\ \hline
NFR-PRIV-7 & Privacy settings shall apply consistently to forum posts, comments, member lists, and search results. \\ \hline
NFR-USE-2 & The user interface shall support responsive layouts for screen widths from 320 px to 1920 px. \\ \hline
NFR-USE-4 & Primary interface actions shall use clear labels and recognizable controls. \\ \hline
Design constraint & Search field names, filters, sort options, component names, and displayed messages follow \texttt{library.md} as UI recommendations; AC/FR/NFR entries are confirmed requirements. \\ \hline
\end{longtable}
\endgroup
\par\medskip

\hypertarget{sec:b9-extension-points}{}
\addcontentsline{toc}{subsubsection}{B9. Extension Points}
\subsubsection*{B9. Extension Points}
\par\medskip

\begingroup
\small
\renewcommand{\baselinestretch}{1.0}\selectfont
\renewcommand{\arraystretch}{1.15}
\setlength{\tabcolsep}{6pt}
\begin{longtable}{|>{\raggedright\arraybackslash}p{0.46\textwidth}|>{\raggedright\arraybackslash}p{0.46\textwidth}|}
\hline
\textbf{Extension Point} & \textbf{Location in Flow of Events} \\ \hline
\endhead
Field-specific search & Basic flow step 5; alternative A1 \\ \hline
Course and academic-scope refinement & Basic flow steps 6-8; alternative A2 \\ \hline
Rating display decision & Basic flow step 9; alternative A3 \\ \hline
Filter reset & Basic flow step 6; alternative A4 \\ \hline
Empty or failed result load & Basic flow step 8; exceptions E1-E2 \\ \hline
Privacy-blocked access & Basic flow step 11; exception E3 \\ \hline
\end{longtable}
\endgroup
\par\medskip

\hypertarget{sec:b10-interface-and-data-notes}{}
\addcontentsline{toc}{subsubsection}{B10. Interface and Data Notes}
\subsubsection*{B10. Interface and Data Notes}
\par\medskip

\begin{itemize}
\item \textbf{Screen and flow path (confirmed):} \texttt{screen\_flow.md} defines \texttt{Home -$>$ Library Group}, \texttt{Library Group -$>$ Exam}, and \texttt{Library Group -$>$ Book}. \texttt{feature\_flow.md} defines \texttt{Library -$>$ Search -$>$ Tag / Title}; \texttt{userstory.md} section 8 defines \texttt{Library -$>$ Search} with Tag and Title.
\item \textbf{Route (not specified):} Route is Not specified in proposals; \texttt{library.md} labels \texttt{/library} as a recommended route only.
\item \textbf{Recommended UI:} \texttt{LibraryPage} / \texttt{LearningResourceLibraryPage}, \texttt{LibraryTabBar} with \textbf{All}, \textbf{Books}, \textbf{Exams}, \textbf{Recommended}, \textbf{Saved}, \textbf{My Storage}, \texttt{ResourceSearchPanel}, \texttt{ResourceKeywordSearchInput}, \textbf{Search}, \textbf{Advanced filters}, \texttt{ActiveFilterChipList}, \textbf{Clear all filters}, and \texttt{ResourceResultGrid}.
\item \textbf{Confirmed search fields:} Title, author, subject tags, and description (\texttt{AC-25.1}, \texttt{FR-25.2}). Recommended field labels include \textbf{Title}, \textbf{Author}, and \textbf{Description contains}, with subject-tag chips such as \texttt{Calculus}, \texttt{Database}, and \texttt{Physics}.
\item \textbf{Confirmed filters and recommended sort:} Filter by format and subject field (\texttt{AC-25.2}, \texttt{FR-25.3}). Recommended format values include PDF, slides, notes, past exam, document, and video; recommended \texttt{ResourceSortSelect} options are \textbf{Relevance}, \textbf{Top rated}, \textbf{Most accessed}, and \textbf{Newest}.
\item \textbf{Exact messages:} “Search title, author, tag, or description.”; “No resources found”; “24 resources found.” as the example \texttt{ResultCountText}; “4.6 (28 reviews)” as the example \texttt{ResourceRatingDisplay}; “Not yet rated.”; “Unable to load resources.” with \textbf{Retry}; “You do not have permission to access this resource.”
\item \textbf{Data fields (confirmed basis plus recommended names):} \texttt{resourceId}, \texttt{title}, \texttt{description}, \texttt{resourceType}, \texttt{author}, \texttt{uploader}, \texttt{subjectTags}, \texttt{courseTags}, \texttt{academicContext}, \texttt{format}, \texttt{privacy}, \texttt{rating}, \texttt{reviewCount}, \texttt{views}, \texttt{accesses}, \texttt{status}, \texttt{publishedAt}, and \texttt{sourceForumPostId} are named in \texttt{library.md}; ratings and review counts display only when available.
\item \textbf{API and codes:} Library search endpoint, statuses, and error codes are Not specified in proposals. \texttt{API.md} has no standalone library search endpoint; \texttt{GET /forum/posts} is a forum-post endpoint, not a library search endpoint, and \texttt{response\_code.md} defines no library-specific error code.
\item \textbf{Notification and privacy rules:} No story-specific notification is stated. Privacy settings apply consistently to search results (\texttt{NFR-PRIV-7}), and resource access follows the owner's privacy setting (\texttt{FR-27.5}, \texttt{NFR-PRIV-5}).
\end{itemize}
\par\medskip

\hypertarget{sec:b11-source-alignment-notes}{}
\addcontentsline{toc}{subsubsection}{B11. Source Alignment Notes}
\subsubsection*{B11. Source Alignment Notes}
\par\medskip

\begin{itemize}
\item This use case is derived from US-25 in \texttt{userstory.md} section 8.1 and from \texttt{FR-25.1} to \texttt{FR-25.5} and \texttt{AC-25.1} to \texttt{AC-25.4} in the proposals.
\item Component names, routes, sort options, example counts, and exact empty/error messages come from \texttt{library.md} as UI/UX recommendations, not as new functional requirements.
\item \texttt{API.md} has no standalone library search endpoint, so no library search route, status code, or error code is invented.
\item Overlap with existing \texttt{UC-06-share-learning-resource.md}: UC-06 makes a published resource searchable and links to the library, while UC-22 covers finding and reviewing those resources; publication details remain in UC-06.
\item No I-xx, E-xx, or G-xx model relationship ID is introduced; A5/A4-style flow labels in this document are specification flow headings, not model IDs.
\end{itemize}
\par\medskip

\par\medskip
\hypertarget{sec:ui-prototype-specimens}{}
\addcontentsline{toc}{subsection}{UI Prototype Specimens}
\subsection*{UI Prototype Specimens}
\par\medskip
\begin{figure}[H]
\centering
\includegraphics[width=0.75\textwidth,height=0.42\textheight,keepaspectratio]{"../../prototype/browse_home_desktop.png"}
\par\vspace{2pt}
{\footnotesize \textbf{UI Prototype:} Browse home - Desktop}
\end{figure}
\par\smallskip
\begin{figure}[H]
\centering
\includegraphics[width=0.75\textwidth,height=0.42\textheight,keepaspectratio]{"../../prototype/browse_computer_science_desktop.png"}
\par\vspace{2pt}
{\footnotesize \textbf{UI Prototype:} Browse: Computer Science - Desktop}
\end{figure}
\par\smallskip

\pagebreak
\hypertarget{sec:uc-23-receive-resource-recommendations}{}
\addcontentsline{toc}{section}{UC-23 --- Receive Resource Recommendations}
\section*{UC-23 --- Receive Resource Recommendations}
\par\medskip

\textbf{Version:} \texttt{1.0}
\par\medskip

\textbf{Performed by:} [Nguyen Le Quoc Huy] | \textbf{Edited by:} [Nguyen Le Quoc Huy]
\par\medskip

\begingroup
\small
\renewcommand{\baselinestretch}{1.0}\selectfont
\renewcommand{\arraystretch}{1.15}
\setlength{\tabcolsep}{6pt}
\begin{longtable}{|>{\raggedright\arraybackslash}p{0.09\textwidth}|>{\raggedright\arraybackslash}p{0.84\textwidth}|}
\hline
\textbf{Field} & \textbf{Value} \\ \hline
\endhead
Use Case ID & UC-23 \\ \hline
Name & Receive Resource Recommendations \\ \hline
Primary Actor & University Student \\ \hline
Supporting Actors & None \\ \hline
Feature Group & Learning Resource Library \\ \hline
Priority & Not explicitly stated \\ \hline
Related User Stories & US-26 \\ \hline
Related Requirements & FR-26.1 to FR-26.4; AC-26.1 to AC-26.3; FR-09.4; FR-09.5; FR-25.5; FR-27.4; NFR-PERF-1; NFR-SCAL-4; NFR-SEC-1; NFR-SEC-6; NFR-PRIV-5; NFR-PRIV-7; NFR-USE-4 \\ \hline
Source Documents & \texttt{docs/requirements/proposals/userstory.md}; \texttt{docs/requirements/proposals/functional\_requirement.md}; \texttt{docs/requirements/proposals/library.md}; \texttt{docs/requirements/proposals/create\_post.md}; \texttt{docs/requirements/proposals/API.md}; \texttt{docs/requirements/proposals/response\_code.md}; \texttt{docs/requirements/proposals/non\_functional\_requirement.md}; \texttt{docs/requirements/proposals/feature\_flow.md}; \texttt{docs/requirements/proposals/screen\_flow.md} \\ \hline
\end{longtable}
\endgroup
\par\medskip

\hypertarget{sec:part-a-use-case-model}{}
\addcontentsline{toc}{subsection}{Part A --- Use Case Model}
\subsection*{Part A --- Use Case Model}
\par\medskip

\hypertarget{sec:a1-focused-use-case-diagram}{}
\addcontentsline{toc}{subsubsection}{A1. Focused Use Case Diagram}
\subsubsection*{A1. Focused Use Case Diagram}
\par\medskip

\begin{center}
\includegraphics[width=0.85\textwidth,height=0.75\textheight,keepaspectratio]{"figures/UC-23-resource-recommendations_1.png"}
\end{center}
\par\medskip

\textbf{Diagram notation.} Person shape = human actor; cloud = supporting system; blue = selected use case; solid line = association; labelled solid arrow = generalization; dashed arrow = include or extend.
\par\medskip

\hypertarget{sec:a2-actor-register}{}
\addcontentsline{toc}{subsubsection}{A2. Actor Register}
\subsubsection*{A2. Actor Register}
\par\medskip

\begingroup
\small
\renewcommand{\baselinestretch}{1.0}\selectfont
\renewcommand{\arraystretch}{1.15}
\setlength{\tabcolsep}{6pt}
\begin{longtable}{|>{\raggedright\arraybackslash}p{0.11\textwidth}|>{\raggedright\arraybackslash}p{0.09\textwidth}|>{\raggedright\arraybackslash}p{0.71\textwidth}|}
\hline
\textbf{Actor} & \textbf{Type} & \textbf{Description} \\ \hline
\endhead
University Student & Primary & Authenticated student who receives resource suggestions based on subjects, courses, interests, and academic context. \\ \hline
\end{longtable}
\endgroup
\par\medskip

\hypertarget{sec:a3-relationships}{}
\addcontentsline{toc}{subsubsection}{A3. Relationships}
\subsubsection*{A3. Relationships}
\par\medskip

\hypertarget{sec:a3-1-associations}{}
\paragraph*{A3.1 Associations}
\par\medskip

\begingroup
\small
\renewcommand{\baselinestretch}{1.0}\selectfont
\renewcommand{\arraystretch}{1.15}
\setlength{\tabcolsep}{6pt}
\begin{longtable}{|>{\raggedright\arraybackslash}p{0.13\textwidth}|>{\raggedright\arraybackslash}p{0.27\textwidth}|>{\raggedright\arraybackslash}p{0.09\textwidth}|>{\raggedright\arraybackslash}p{0.41\textwidth}|}
\hline
\textbf{Actor} & \textbf{Use Case} & \textbf{Type} & \textbf{Description} \\ \hline
\endhead
University Student & UC-23 Receive Resource Recommendations & Primary & Views and saves recommended books and learning resources. \\ \hline
\end{longtable}
\endgroup
\par\medskip

\hypertarget{sec:a3-2-generalization-and-specialization}{}
\paragraph*{A3.2 Generalization and Specialization}
\par\medskip

\begingroup
\small
\renewcommand{\baselinestretch}{1.0}\selectfont
\renewcommand{\arraystretch}{1.15}
\setlength{\tabcolsep}{6pt}
\begin{longtable}{|>{\raggedright\arraybackslash}p{0.09\textwidth}|>{\raggedright\arraybackslash}p{0.09\textwidth}|>{\raggedright\arraybackslash}p{0.19\textwidth}|>{\raggedright\arraybackslash}p{0.54\textwidth}|}
\hline
\textbf{Child} & \textbf{Parent} & \textbf{Type} & \textbf{Inherited Behavior or Association} \\ \hline
\endhead
None & None & No generalization or specialization & The proposals describe new users as a recommendation condition, not as a distinct actor for US-26. \\ \hline
\end{longtable}
\endgroup
\par\medskip

\hypertarget{sec:a3-3-include}{}
\paragraph*{A3.3 Include}
\par\medskip

\begingroup
\small
\renewcommand{\baselinestretch}{1.0}\selectfont
\renewcommand{\arraystretch}{1.15}
\setlength{\tabcolsep}{6pt}
\begin{longtable}{|>{\raggedright\arraybackslash}p{0.23\textwidth}|>{\raggedright\arraybackslash}p{0.30\textwidth}|>{\raggedright\arraybackslash}p{0.40\textwidth}|}
\hline
\textbf{Base Use Case} & \textbf{Included Use Case} & \textbf{Meaning} \\ \hline
\endhead
None & None & No include relationship & The proposals do not name an included use case for US-26. \\ \hline
\end{longtable}
\endgroup
\par\medskip

\hypertarget{sec:a3-4-extend}{}
\paragraph*{A3.4 Extend}
\par\medskip

\begingroup
\small
\renewcommand{\baselinestretch}{1.0}\selectfont
\renewcommand{\arraystretch}{1.15}
\setlength{\tabcolsep}{6pt}
\begin{longtable}{|>{\raggedright\arraybackslash}p{0.15\textwidth}|>{\raggedright\arraybackslash}p{0.11\textwidth}|>{\raggedright\arraybackslash}p{0.18\textwidth}|>{\raggedright\arraybackslash}p{0.47\textwidth}|}
\hline
\textbf{Extending Use Case} & \textbf{Base Use Case} & \textbf{Condition} & \textbf{Added Behavior} \\ \hline
\endhead
None & None & No extend relationship & The proposals do not name an extending use case for US-26. \\ \hline
\end{longtable}
\endgroup
\par\medskip

\hypertarget{sec:a4-model-notes}{}
\addcontentsline{toc}{subsubsection}{A4. Model Notes}
\subsubsection*{A4. Model Notes}
\par\medskip

\begin{enumerate}
\item The scope is personalized and cold-start resource suggestions, their display, and saving recommended resources to a personal list.
\item Searching by query, publishing resources, and deleting published resources are outside this use case.
\item Dismissal and refresh controls are library.md UI recommendations; they are not modeled as new extending use cases.
\item No distinct included, extending, or specialized use case is named for US-26, so no I-xx, E-xx, or G-xx relationship ID is introduced.
\item The exact API endpoint, route, recommendation count, and priority are not specified in proposals.
\end{enumerate}
\par\medskip

\hypertarget{sec:part-b-use-case-specification}{}
\addcontentsline{toc}{subsection}{Part B --- Use Case Specification}
\subsection*{Part B --- Use Case Specification}
\par\medskip

\hypertarget{sec:b1-brief-description}{}
\addcontentsline{toc}{subsubsection}{B1. Brief Description}
\subsubsection*{B1. Brief Description}
\par\medskip

A university student receives recommendations for books and learning resources so that the student can discover useful materials related to subjects. The system recommends resources related to the student's subjects, courses, and interests, displays them on the home or library page, and allows the student to save recommended resources to a personal list.
\par\medskip

For new users without historical activity, the system provides curated foundational resources. The library project frames discovery around \textbf{For you}, \textbf{Top reading}, \textbf{Course progress}, and \textbf{Starter set}, with user-relevance and course-progress recommendations drawn from the feature flows.
\par\medskip

\hypertarget{sec:b2-preconditions}{}
\addcontentsline{toc}{subsubsection}{B2. Preconditions}
\subsubsection*{B2. Preconditions}
\par\medskip

\begin{enumerate}
\item The student is authenticated before accessing protected user resources (\texttt{NFR-SEC-1}).
\item The Learning Resource Library is reachable through \texttt{Home -$>$ Library Group} in \texttt{screen\_flow.md}.
\item The student's preferences, courses, subjects, and interests are available where the student has provided them (\texttt{FR-09.1}, \texttt{FR-09.2}, \texttt{FR-09.5}; \texttt{AC-09.1}, \texttt{AC-09.2}, \texttt{AC-09.5}).
\item Published resources are available for recommendation after publication (\texttt{FR-27.4}, \texttt{AC-27.1}).
\item For a new user, limited historical activity is treated as the cold-start condition (\texttt{FR-26.4}, \texttt{AC-26.3}).
\end{enumerate}
\par\medskip

\hypertarget{sec:b3-postconditions}{}
\addcontentsline{toc}{subsubsection}{B3. Postconditions}
\subsubsection*{B3. Postconditions}
\par\medskip

\begin{enumerate}
\item \textbf{Success:} Recommended resources are displayed on the home or library page.
\item \textbf{Success:} Suggestions reflect the student's preferences and courses, or curated foundational resources for a new user.
\item \textbf{Success:} A recommended resource selected for saving is added to the student's personal list.
\item \textbf{Failure:} No recommendation is presented as personalized when the system has no basis, and the specified cold-start or load-error feedback is shown.
\end{enumerate}
\par\medskip

\hypertarget{sec:b4-main-success-scenario-basic-flow}{}
\addcontentsline{toc}{subsubsection}{B4. Main Success Scenario (Basic Flow)}
\subsubsection*{B4. Main Success Scenario (Basic Flow)}
\par\medskip

\begin{enumerate}
\item The student opens the library recommendation area or the home-page recommendation area.
\item The system displays \texttt{RecommendationPanel} with heading \textbf{Recommended for your studies}.
\item The student selects a recommendation segment such as \textbf{For you}, \textbf{Top reading}, \textbf{Course progress}, or \textbf{Starter set}.
\item The system recommends books and learning resources related to the student's subjects, courses, and interests (\texttt{FR-26.1}, \texttt{AC-26.1}).
\item The system displays each recommendation with \texttt{RecommendationReasonText}, for example “Recommended because you are enrolled in CS101.”
\item The student reviews the recommendation card with title, type, tags, rating, and access information.
\item The student selects \textbf{Save} through \texttt{SaveRecommendationButton}.
\item The system saves the recommended resource to the student's personal list (\texttt{FR-26.3}, \texttt{AC-26.2}).
\item The saved item becomes available through \texttt{SavedResourceList}.
\end{enumerate}
\par\medskip

\hypertarget{sec:b5-alternative-flows}{}
\addcontentsline{toc}{subsubsection}{B5. Alternative Flows}
\subsubsection*{B5. Alternative Flows}
\par\medskip

\hypertarget{sec:a1-home-page-display}{}
\paragraph*{A1. Home-page display}
\par\medskip

The recommendation set is shown on the student's home page instead of the library page (\texttt{FR-26.2}). The same save action adds the resource to the personal list.
\par\medskip

\hypertarget{sec:a2-top-reading-or-course-progress}{}
\paragraph*{A2. Top reading or course progress}
\par\medskip

The student selects \textbf{Top reading} or \textbf{Course progress}. The system displays popular or course-connected resources with rating and access information; these are library.md recommendation segments derived from \texttt{feature\_flow.md} and the course-progress flow.
\par\medskip

\hypertarget{sec:a3-cold-start-starter-resources}{}
\paragraph*{A3. Cold-start starter resources}
\par\medskip

The student is new or has little historical activity. The system displays \texttt{CuratedStarterResourceList} with foundational resources based on selected subjects (\texttt{FR-26.4}, \texttt{AC-26.3}).
\par\medskip

\hypertarget{sec:a4-dismiss-or-refresh-suggestions}{}
\paragraph*{A4. Dismiss or refresh suggestions}
\par\medskip

The student selects \textbf{Not relevant} to hide an item or \textbf{Refresh suggestions} to request a new set. These controls are explicit library.md recommendations and are not confirmed FRs for US-26.
\par\medskip

\hypertarget{sec:b6-exception-flows}{}
\addcontentsline{toc}{subsubsection}{B6. Exception Flows}
\subsubsection*{B6. Exception Flows}
\par\medskip

\hypertarget{sec:e1-recommendation-load-failure}{}
\paragraph*{E1. Recommendation load failure}
\par\medskip

The recommendation set cannot be loaded. The system shows “Unable to load resources.” with \textbf{Retry} through \texttt{LibraryErrorBanner}.
\par\medskip

\hypertarget{sec:e2-recommendation-reason-unavailable}{}
\paragraph*{E2. Recommendation reason unavailable}
\par\medskip

No explanation basis is available for a suggested resource. The system omits \texttt{RecommendationReasonText} rather than displaying an invented reason.
\par\medskip

\hypertarget{sec:e3-save-failure}{}
\paragraph*{E3. Save failure}
\par\medskip

The save action does not complete. The system keeps the recommendation available for another attempt; the exact save-failure message is Not specified in proposals.
\par\medskip

\hypertarget{sec:b7-business-rules}{}
\addcontentsline{toc}{subsubsection}{B7. Business Rules}
\subsubsection*{B7. Business Rules}
\par\medskip

\begin{enumerate}
\item The system recommends books or learning resources related to the user's subjects, courses, and interests (\texttt{FR-26.1}, \texttt{AC-26.1}).
\item The system displays recommended resources on the user's home or library page (\texttt{FR-26.2}).
\item Users can save recommended resources to a personal list (\texttt{FR-26.3}, \texttt{AC-26.2}).
\item New users without historical activity receive curated foundational resources (\texttt{FR-26.4}, \texttt{AC-26.3}).
\item Learning preferences are used to recommend materials (\texttt{FR-09.4}), and preferences can be updated at any time (\texttt{FR-09.5}, \texttt{AC-09.5}).
\item The library displays resources related to the user's university, major, and courses (\texttt{FR-25.5}, \texttt{AC-25.4}).
\item Access to recommended resources follows owner privacy settings (\texttt{FR-27.5}, \texttt{NFR-PRIV-5}, \texttt{NFR-PRIV-7}).
\end{enumerate}
\par\medskip

\hypertarget{sec:b8-special-requirements}{}
\addcontentsline{toc}{subsubsection}{B8. Special Requirements}
\subsubsection*{B8. Special Requirements}
\par\medskip

\begingroup
\small
\renewcommand{\baselinestretch}{1.0}\selectfont
\renewcommand{\arraystretch}{1.15}
\setlength{\tabcolsep}{6pt}
\begin{longtable}{|>{\raggedright\arraybackslash}p{0.10\textwidth}|>{\raggedright\arraybackslash}p{0.83\textwidth}|}
\hline
\textbf{ID} & \textbf{Requirement} \\ \hline
\endhead
NFR-PERF-1 & The system shall provide responsive user interactions with a response time of ≤ 2 seconds for standard user actions. \\ \hline
NFR-SCAL-4 & The system architecture shall support at least 100,000 uploaded documents and learning resources. \\ \hline
NFR-SEC-1 & The system shall require authentication before accessing protected user resources. \\ \hline
NFR-SEC-6 & The system shall prevent unauthorized access through server-side authorization checks. \\ \hline
NFR-PRIV-5 & Users shall be able to control access to documents and learning resources they share. \\ \hline
NFR-PRIV-7 & Privacy settings shall apply consistently to forum posts, comments, member lists, and search results. \\ \hline
NFR-USE-4 & Primary interface actions shall use clear labels and recognizable controls. \\ \hline
Design constraint & Recommendation segments, reason text, and dismissal/refresh controls follow \texttt{library.md} as UI recommendations; FR/AC/NFR entries are confirmed requirements. \\ \hline
\end{longtable}
\endgroup
\par\medskip

\hypertarget{sec:b9-extension-points}{}
\addcontentsline{toc}{subsubsection}{B9. Extension Points}
\subsubsection*{B9. Extension Points}
\par\medskip

\begingroup
\small
\renewcommand{\baselinestretch}{1.0}\selectfont
\renewcommand{\arraystretch}{1.15}
\setlength{\tabcolsep}{6pt}
\begin{longtable}{|>{\raggedright\arraybackslash}p{0.37\textwidth}|>{\raggedright\arraybackslash}p{0.56\textwidth}|}
\hline
\textbf{Extension Point} & \textbf{Location in Flow of Events} \\ \hline
\endhead
Home or library placement & Basic flow step 1; alternative A1 \\ \hline
Segment selection & Basic flow step 3; alternative A2 \\ \hline
Cold-start handling & Before basic flow step 4; alternative A3 \\ \hline
Dismiss or refresh & After basic flow step 5; alternative A4 \\ \hline
Recommendation load failure & Basic flow steps 2-4; exception E1 \\ \hline
Save failure & Basic flow steps 7-8; exception E3 \\ \hline
\end{longtable}
\endgroup
\par\medskip

\hypertarget{sec:b10-interface-and-data-notes}{}
\addcontentsline{toc}{subsubsection}{B10. Interface and Data Notes}
\subsubsection*{B10. Interface and Data Notes}
\par\medskip

\begin{itemize}
\item \textbf{Screen and flow path (confirmed):} \texttt{screen\_flow.md} defines \texttt{Home -$>$ Library Group}. \texttt{feature\_flow.md} defines \texttt{Library -$>$ Recommend -$>$ Top reading / User relevance}; \texttt{userstory.md} section 8 also shows \texttt{Recommend} with Top reading and User relevance, and \texttt{library.md} includes course progress.
\item \textbf{Route (not specified):} Route is Not specified in proposals; \texttt{library.md} presents \texttt{/library/recommended} and \texttt{/library/saved} as suggested states only.
\item \textbf{Recommended UI:} \texttt{RecommendationPanel} heading \textbf{Recommended for your studies}, \texttt{RecommendationSegmentControl} options \textbf{For you}, \textbf{Top reading}, \textbf{Course progress}, \textbf{Starter set}, \texttt{UserRelevanceResourceList}, \texttt{TopReadingResourceList}, \texttt{CourseProgressResourceList}, \texttt{CuratedStarterResourceList}, and \texttt{SavedResourceList}.
\item \textbf{Confirmed actions and data:} Display recommendations (\texttt{FR-26.2}, \texttt{AC-26.1}) and save them to a personal list (\texttt{FR-26.3}, \texttt{AC-26.2}). Recommended fields include \texttt{recommendationReason}, \texttt{title}, \texttt{resourceType}, \texttt{subjectTags}, \texttt{courseTags}, \texttt{rating}, \texttt{reviewCount}, \texttt{accesses}, \texttt{status}, and \texttt{publishedAt}.
\item \textbf{Exact messages and labels:} \textbf{Recommended for your studies}; \textbf{For you}; \textbf{Top reading}; \textbf{Course progress}; \textbf{Starter set}; “Recommended because you are enrolled in CS101.”; \textbf{Save}; \textbf{Not relevant}; \textbf{Refresh suggestions}; “Unable to load resources.” with \textbf{Retry}.
\item \textbf{API and codes:} Library recommendation endpoint, statuses, and error codes are Not specified in proposals. \texttt{API.md} has no standalone library recommendation endpoint, and \texttt{response\_code.md} defines no library-specific error code.
\item \textbf{Recommendation count and limits:} The number of recommendations, refresh interval, and numeric cold-start threshold are Not specified in proposals.
\item \textbf{Notification and privacy rules:} No story-specific notification is stated. Recommendations must respect resource privacy settings (\texttt{FR-27.5}, \texttt{NFR-PRIV-5}, \texttt{NFR-PRIV-7}).
\end{itemize}
\par\medskip

\hypertarget{sec:b11-source-alignment-notes}{}
\addcontentsline{toc}{subsubsection}{B11. Source Alignment Notes}
\subsubsection*{B11. Source Alignment Notes}
\par\medskip

\begin{itemize}
\item This use case is derived from US-26 in \texttt{userstory.md} section 8.2 and from \texttt{FR-26.1} to \texttt{FR-26.4} and \texttt{AC-26.1} to \texttt{AC-26.3} in the proposals.
\item Preference and relevance dependencies come from \texttt{FR-09.4}, \texttt{FR-09.5}, and \texttt{FR-25.5}; publication dependency comes from \texttt{FR-27.4}.
\item Recommendation segments, reason text, dismissal, refresh, and exact UI labels come from \texttt{library.md} as UI/UX recommendations; dismissal is analogous to \texttt{FR-24.4} but is not a confirmed US-26 FR.
\item \texttt{API.md} has no standalone library recommendation endpoint, so no library recommendation route, status code, or error code is invented.
\item Overlap with existing \texttt{UC-06-share-learning-resource.md}: UC-06 publishes resources into the library; UC-23 surfaces published resources as recommendations and saves them to a personal list.
\item No I-xx, E-xx, or G-xx model relationship ID is introduced; flow headings such as A3 and E1 are specification flow labels only.
\end{itemize}
\par\medskip

\par\medskip
\hypertarget{sec:ui-prototype-specimens}{}
\addcontentsline{toc}{subsection}{UI Prototype Specimens}
\subsection*{UI Prototype Specimens}
\par\medskip
\begin{figure}[H]
\centering
\includegraphics[width=0.75\textwidth,height=0.42\textheight,keepaspectratio]{"../../prototype/feed_algorithms_desktop.png"}
\par\vspace{2pt}
{\footnotesize \textbf{UI Prototype:} Feed: Algorithms - Desktop}
\end{figure}
\par\smallskip

\pagebreak
\hypertarget{sec:uc-24-view-my-storage}{}
\addcontentsline{toc}{section}{UC-24 --- View My Storage}
\section*{UC-24 --- View My Storage}
\par\medskip

\textbf{Version:} \texttt{1.0}
\par\medskip

\textbf{Performed by:} [Nguyen Le Quoc Huy] | \textbf{Edited by:} [Nguyen Le Quoc Huy]
\par\medskip

\begingroup
\small
\renewcommand{\baselinestretch}{1.0}\selectfont
\renewcommand{\arraystretch}{1.15}
\setlength{\tabcolsep}{6pt}
\begin{longtable}{|>{\raggedright\arraybackslash}p{0.09\textwidth}|>{\raggedright\arraybackslash}p{0.84\textwidth}|}
\hline
\textbf{Field} & \textbf{Value} \\ \hline
\endhead
Use Case ID & UC-24 \\ \hline
Name & View My Storage \\ \hline
Primary Actor & University Student \\ \hline
Supporting Actors & None \\ \hline
Feature Group & Learning Resource Library \\ \hline
Priority & Not explicitly stated \\ \hline
Related User Stories & US-28 \\ \hline
Related Requirements & FR-28.1 to FR-28.3; AC-28.1 to AC-28.3; FR-27.1; FR-27.5; NFR-PERF-1; NFR-SCAL-4; NFR-SEC-1; NFR-SEC-6; NFR-PRIV-5; NFR-USE-4 \\ \hline
Source Documents & \texttt{docs/requirements/proposals/userstory.md}; \texttt{docs/requirements/proposals/functional\_requirement.md}; \texttt{docs/requirements/proposals/library.md}; \texttt{docs/requirements/proposals/create\_post.md}; \texttt{docs/requirements/proposals/API.md}; \texttt{docs/requirements/proposals/response\_code.md}; \texttt{docs/requirements/proposals/non\_functional\_requirement.md}; \texttt{docs/requirements/proposals/feature\_flow.md}; \texttt{docs/requirements/proposals/screen\_flow.md} \\ \hline
\end{longtable}
\endgroup
\par\medskip

\hypertarget{sec:part-a-use-case-model}{}
\addcontentsline{toc}{subsection}{Part A --- Use Case Model}
\subsection*{Part A --- Use Case Model}
\par\medskip

\hypertarget{sec:a1-focused-use-case-diagram}{}
\addcontentsline{toc}{subsubsection}{A1. Focused Use Case Diagram}
\subsubsection*{A1. Focused Use Case Diagram}
\par\medskip

\begin{center}
\includegraphics[width=0.85\textwidth,height=0.75\textheight,keepaspectratio]{"figures/UC-24-view-my-storage_1.png"}
\end{center}
\par\medskip

\textbf{Diagram notation.} Person shape = human actor; cloud = supporting system; blue = selected use case; solid line = association; labelled solid arrow = generalization; dashed arrow = include or extend.
\par\medskip

\hypertarget{sec:a2-actor-register}{}
\addcontentsline{toc}{subsubsection}{A2. Actor Register}
\subsubsection*{A2. Actor Register}
\par\medskip

\begingroup
\small
\renewcommand{\baselinestretch}{1.0}\selectfont
\renewcommand{\arraystretch}{1.15}
\setlength{\tabcolsep}{6pt}
\begin{longtable}{|>{\raggedright\arraybackslash}p{0.13\textwidth}|>{\raggedright\arraybackslash}p{0.09\textwidth}|>{\raggedright\arraybackslash}p{0.69\textwidth}|}
\hline
\textbf{Actor} & \textbf{Type} & \textbf{Description} \\ \hline
\endhead
University Student & Primary & Authenticated resource owner who reviews published resources and their view and access counts. \\ \hline
\end{longtable}
\endgroup
\par\medskip

\hypertarget{sec:a3-relationships}{}
\addcontentsline{toc}{subsubsection}{A3. Relationships}
\subsubsection*{A3. Relationships}
\par\medskip

\hypertarget{sec:a3-1-associations}{}
\paragraph*{A3.1 Associations}
\par\medskip

\begingroup
\small
\renewcommand{\baselinestretch}{1.0}\selectfont
\renewcommand{\arraystretch}{1.15}
\setlength{\tabcolsep}{6pt}
\begin{longtable}{|>{\raggedright\arraybackslash}p{0.16\textwidth}|>{\raggedright\arraybackslash}p{0.18\textwidth}|>{\raggedright\arraybackslash}p{0.10\textwidth}|>{\raggedright\arraybackslash}p{0.47\textwidth}|}
\hline
\textbf{Actor} & \textbf{Use Case} & \textbf{Type} & \textbf{Description} \\ \hline
\endhead
University Student & UC-24 View My Storage & Primary & Opens a personal storage view of published resources. \\ \hline
\end{longtable}
\endgroup
\par\medskip

\hypertarget{sec:a3-2-generalization-and-specialization}{}
\paragraph*{A3.2 Generalization and Specialization}
\par\medskip

\begingroup
\small
\renewcommand{\baselinestretch}{1.0}\selectfont
\renewcommand{\arraystretch}{1.15}
\setlength{\tabcolsep}{6pt}
\begin{longtable}{|>{\raggedright\arraybackslash}p{0.09\textwidth}|>{\raggedright\arraybackslash}p{0.09\textwidth}|>{\raggedright\arraybackslash}p{0.26\textwidth}|>{\raggedright\arraybackslash}p{0.47\textwidth}|}
\hline
\textbf{Child} & \textbf{Parent} & \textbf{Type} & \textbf{Inherited Behavior or Association} \\ \hline
\endhead
None & None & No generalization or specialization & The proposals name only the University Student actor for US-28. \\ \hline
\end{longtable}
\endgroup
\par\medskip

\hypertarget{sec:a3-3-include}{}
\paragraph*{A3.3 Include}
\par\medskip

\begingroup
\small
\renewcommand{\baselinestretch}{1.0}\selectfont
\renewcommand{\arraystretch}{1.15}
\setlength{\tabcolsep}{6pt}
\begin{longtable}{|>{\raggedright\arraybackslash}p{0.23\textwidth}|>{\raggedright\arraybackslash}p{0.30\textwidth}|>{\raggedright\arraybackslash}p{0.40\textwidth}|}
\hline
\textbf{Base Use Case} & \textbf{Included Use Case} & \textbf{Meaning} \\ \hline
\endhead
None & None & No include relationship & The proposals do not name an included use case for US-28. \\ \hline
\end{longtable}
\endgroup
\par\medskip

\hypertarget{sec:a3-4-extend}{}
\paragraph*{A3.4 Extend}
\par\medskip

\begingroup
\small
\renewcommand{\baselinestretch}{1.0}\selectfont
\renewcommand{\arraystretch}{1.15}
\setlength{\tabcolsep}{6pt}
\begin{longtable}{|>{\raggedright\arraybackslash}p{0.15\textwidth}|>{\raggedright\arraybackslash}p{0.11\textwidth}|>{\raggedright\arraybackslash}p{0.18\textwidth}|>{\raggedright\arraybackslash}p{0.47\textwidth}|}
\hline
\textbf{Extending Use Case} & \textbf{Base Use Case} & \textbf{Condition} & \textbf{Added Behavior} \\ \hline
\endhead
None & None & No extend relationship & The proposals do not name an extending use case for US-28. \\ \hline
\end{longtable}
\endgroup
\par\medskip

\hypertarget{sec:a4-model-notes}{}
\addcontentsline{toc}{subsubsection}{A4. Model Notes}
\subsubsection*{A4. Model Notes}
\par\medskip

\begin{enumerate}
\item The scope is the owner's personal storage view, status and publication-date display, and view/access statistics.
\item Publishing resources is outside this use case; deletion is handled in US-29 and existing UC-06 covers the publication flow.
\item The status values shown in the storage table are exactly \texttt{Published}, \texttt{Restricted}, and \texttt{Unavailable}.
\item No distinct included, extending, or specialized use case is named for US-28, so no I-xx, E-xx, or G-xx relationship ID is introduced.
\item The exact API endpoint, route, pagination limit, and priority are not specified in proposals.
\end{enumerate}
\par\medskip

\hypertarget{sec:part-b-use-case-specification}{}
\addcontentsline{toc}{subsection}{Part B --- Use Case Specification}
\subsection*{Part B --- Use Case Specification}
\par\medskip

\hypertarget{sec:b1-brief-description}{}
\addcontentsline{toc}{subsubsection}{B1. Brief Description}
\subsubsection*{B1. Brief Description}
\par\medskip

A university student views personal storage so that the student can track how many views and accesses published resources receive. The system provides a personal storage view for published resources, displays view and access counts for each resource, and displays publication status and date.
\par\medskip

\texttt{library.md} recommends presenting this as \textbf{My Storage}, with summary totals and a table of owned resources. These UI details are recommendations; the confirmed requirements are \texttt{AC-28.1} to \texttt{AC-28.3} and \texttt{FR-28.1} to \texttt{FR-28.3}.
\par\medskip

\hypertarget{sec:b2-preconditions}{}
\addcontentsline{toc}{subsubsection}{B2. Preconditions}
\subsubsection*{B2. Preconditions}
\par\medskip

\begin{enumerate}
\item The student is authenticated before accessing protected user resources (\texttt{NFR-SEC-1}).
\item The Learning Resource Library is reachable through \texttt{Home -$>$ Library Group} in \texttt{screen\_flow.md}.
\item The student has published resources through the library or material-post path (\texttt{FR-27.1}, \texttt{FR-15.6}, \texttt{FR-27.4}).
\item View, access, status, and publication-date data are available for the owned resources (\texttt{FR-28.2}, \texttt{FR-28.3}).
\item Authorization checks restrict storage data to the resource owner (\texttt{NFR-SEC-6}).
\end{enumerate}
\par\medskip

\hypertarget{sec:b3-postconditions}{}
\addcontentsline{toc}{subsubsection}{B3. Postconditions}
\subsubsection*{B3. Postconditions}
\par\medskip

\begin{enumerate}
\item \textbf{Success:} The student's published resources are listed in a personal storage view.
\item \textbf{Success:} Each listed resource shows view count, access count, publication status, and publication date.
\item \textbf{Success:} The student can open a resource detail view or proceed to the delete action where available.
\item \textbf{Failure:} No other student's storage data is exposed, and the specified empty-state or load-error feedback is shown.
\end{enumerate}
\par\medskip

\hypertarget{sec:b4-main-success-scenario-basic-flow}{}
\addcontentsline{toc}{subsubsection}{B4. Main Success Scenario (Basic Flow)}
\subsubsection*{B4. Main Success Scenario (Basic Flow)}
\par\medskip

\begin{enumerate}
\item The student opens the library tab bar and selects \textbf{My Storage} through \texttt{MyStorageTab}.
\item The system displays \texttt{MyStoragePanel} with heading \textbf{My Storage}.
\item The system displays \texttt{MyStorageSummaryStrip} with total published resources, total views, and total accesses/downloads.
\item The system displays \texttt{MyPublishedResourcesTable} with one row per published resource.
\item Each row shows the resource title and type through \texttt{MyResourceTitleCell}.
\item Each row shows a \texttt{status} badge with exactly \texttt{Published}, \texttt{Restricted}, or \texttt{Unavailable}.
\item Each row shows the publication date and time through \texttt{MyResourcePublishedDateCell} (\texttt{FR-28.3}, \texttt{AC-28.3}).
\item Each row shows the number of views through \texttt{MyResourceViewsCell} and the number of downloads/accesses through \texttt{MyResourceAccessesCell} (\texttt{FR-28.2}, \texttt{AC-28.2}).
\item Each row shows the privacy setting through \texttt{MyResourcePrivacyCell}.
\item The student selects \textbf{View} through \texttt{MyResourceOpenButton} to open the resource detail view, or remains in My Storage.
\end{enumerate}
\par\medskip

\hypertarget{sec:b5-alternative-flows}{}
\addcontentsline{toc}{subsubsection}{B5. Alternative Flows}
\subsubsection*{B5. Alternative Flows}
\par\medskip

\hypertarget{sec:a1-no-published-resources}{}
\paragraph*{A1. No published resources}
\par\medskip

The student has no published resources. The system displays the exact empty-state message “You have not published any resources yet.”
\par\medskip

\hypertarget{sec:a2-restricted-or-unavailable-resource}{}
\paragraph*{A2. Restricted or unavailable resource}
\par\medskip

A listed resource has status \texttt{Restricted} or \texttt{Unavailable}. The system still shows the owner's status badge and statistics, while resource access follows the owner's privacy setting (\texttt{FR-27.5}, \texttt{NFR-PRIV-5}).
\par\medskip

\hypertarget{sec:a3-open-resource-detail}{}
\paragraph*{A3. Open resource detail}
\par\medskip

The student selects \textbf{View}. The system opens the resource detail page with full metadata, preview where supported, rating, access information, views, and accesses.
\par\medskip

\hypertarget{sec:a4-retry-after-an-error}{}
\paragraph*{A4. Retry after an error}
\par\medskip

The student selects \textbf{Retry} through \texttt{StorageRetryButton}. The system reloads the storage data.
\par\medskip

\hypertarget{sec:b6-exception-flows}{}
\addcontentsline{toc}{subsubsection}{B6. Exception Flows}
\subsubsection*{B6. Exception Flows}
\par\medskip

\hypertarget{sec:e1-storage-data-load-failure}{}
\paragraph*{E1. Storage data load failure}
\par\medskip

The storage view cannot be loaded. The system shows “Unable to load resources.” with \textbf{Retry} through \texttt{LibraryErrorBanner}.
\par\medskip

\hypertarget{sec:e2-statistics-unavailable}{}
\paragraph*{E2. Statistics unavailable}
\par\medskip

View or access counts are unavailable for a resource. The system does not present an invented number; the exact missing-statistics message is Not specified in proposals.
\par\medskip

\hypertarget{sec:e3-unauthorized-storage-access}{}
\paragraph*{E3. Unauthorized storage access}
\par\medskip

A user who is not the owner attempts to open another student's storage view. The system denies access through server-side authorization checks (\texttt{NFR-SEC-6}); the exact storage-denial message is Not specified in proposals.
\par\medskip

\hypertarget{sec:b7-business-rules}{}
\addcontentsline{toc}{subsubsection}{B7. Business Rules}
\subsubsection*{B7. Business Rules}
\par\medskip

\begin{enumerate}
\item The system provides a personal storage view for published resources (\texttt{FR-28.1}, \texttt{AC-28.1}).
\item The system displays the number of views and accesses for each published resource (\texttt{FR-28.2}, \texttt{AC-28.2}).
\item The system displays the publication status and date of each resource (\texttt{FR-28.3}, \texttt{AC-28.3}).
\item Storage status values are exactly \texttt{Published}, \texttt{Restricted}, and \texttt{Unavailable}.
\item Privacy settings control access to documents and learning resources the student shared (\texttt{FR-27.5}, \texttt{NFR-PRIV-5}).
\item Only the owner may view personal storage data (\texttt{NFR-SEC-1}, \texttt{NFR-SEC-6}).
\end{enumerate}
\par\medskip

\hypertarget{sec:b8-special-requirements}{}
\addcontentsline{toc}{subsubsection}{B8. Special Requirements}
\subsubsection*{B8. Special Requirements}
\par\medskip

\begingroup
\small
\renewcommand{\baselinestretch}{1.0}\selectfont
\renewcommand{\arraystretch}{1.15}
\setlength{\tabcolsep}{6pt}
\begin{longtable}{|>{\raggedright\arraybackslash}p{0.10\textwidth}|>{\raggedright\arraybackslash}p{0.83\textwidth}|}
\hline
\textbf{ID} & \textbf{Requirement} \\ \hline
\endhead
NFR-PERF-1 & The system shall provide responsive user interactions with a response time of ≤ 2 seconds for standard user actions. \\ \hline
NFR-SCAL-4 & The system architecture shall support at least 100,000 uploaded documents and learning resources. \\ \hline
NFR-SEC-1 & The system shall require authentication before accessing protected user resources. \\ \hline
NFR-SEC-6 & The system shall prevent unauthorized access through server-side authorization checks. \\ \hline
NFR-PRIV-5 & Users shall be able to control access to documents and learning resources they share. \\ \hline
NFR-USE-4 & Primary interface actions shall use clear labels and recognizable controls. \\ \hline
Design constraint & Status values, table cells, summary cards, empty state, and labels follow \texttt{library.md} as UI recommendations; FR/AC/NFR entries are confirmed requirements. \\ \hline
\end{longtable}
\endgroup
\par\medskip

\hypertarget{sec:b9-extension-points}{}
\addcontentsline{toc}{subsubsection}{B9. Extension Points}
\subsubsection*{B9. Extension Points}
\par\medskip

\begingroup
\small
\renewcommand{\baselinestretch}{1.0}\selectfont
\renewcommand{\arraystretch}{1.15}
\setlength{\tabcolsep}{6pt}
\begin{longtable}{|>{\raggedright\arraybackslash}p{0.41\textwidth}|>{\raggedright\arraybackslash}p{0.52\textwidth}|}
\hline
\textbf{Extension Point} & \textbf{Location in Flow of Events} \\ \hline
\endhead
Empty storage & Before basic flow step 4; alternative A1 \\ \hline
Restricted or unavailable status & Basic flow step 6; alternative A2 \\ \hline
Resource detail access & Basic flow step 10; alternative A3 \\ \hline
Storage retry & After exception E1; alternative A4 \\ \hline
Storage load failure & Basic flow steps 2-4; exception E1 \\ \hline
Unauthorized access & Before basic flow step 2; exception E3 \\ \hline
\end{longtable}
\endgroup
\par\medskip

\hypertarget{sec:b10-interface-and-data-notes}{}
\addcontentsline{toc}{subsubsection}{B10. Interface and Data Notes}
\subsubsection*{B10. Interface and Data Notes}
\par\medskip

\begin{itemize}
\item \textbf{Screen and flow path (confirmed):} \texttt{screen\_flow.md} defines \texttt{Home -$>$ Library Group}. \texttt{feature\_flow.md} defines \texttt{Library -$>$ Material -$>$ Upload / Download}; \texttt{userstory.md} section 8.4 defines the personal storage requirement.
\item \textbf{Route (not specified):} Route is Not specified in proposals; \texttt{library.md} presents \texttt{/library/my-storage} as a suggested state only.
\item \textbf{Recommended UI:} \texttt{MyStorageTab} label \textbf{My Storage}, \texttt{MyStoragePanel} heading \textbf{My Storage}, \texttt{MyStorageSummaryStrip}, \texttt{MyPublishedResourcesTable}, \texttt{MyResourceStatusCell}, \texttt{MyResourcePublishedDateCell}, \texttt{MyResourceViewsCell}, \texttt{MyResourceAccessesCell}, \texttt{MyResourcePrivacyCell}, \texttt{MyResourceOpenButton}, \texttt{DeleteMyResourceButton}, and \texttt{StorageRetryButton}.
\item \textbf{Confirmed statistics and status:} \texttt{FR-28.2} / \texttt{AC-28.2} require views and accesses; \texttt{FR-28.3} / \texttt{AC-28.3} require publication status and date. The exact status values are \texttt{Published}, \texttt{Restricted}, and \texttt{Unavailable}.
\item \textbf{Exact messages and labels:} \textbf{My Storage}; \textbf{View}; \textbf{Delete}; \textbf{Retry}; “You have not published any resources yet.”; “Unable to load resources.” with \textbf{Retry}. No exact success message is specified for opening storage.
\item \textbf{Data fields (confirmed basis plus recommended names):} \texttt{resourceId}, \texttt{title}, \texttt{resourceType}, \texttt{privacy}, \texttt{views}, \texttt{accesses}, \texttt{status}, \texttt{publishedAt}, \texttt{rating}, \texttt{reviewCount}, and \texttt{sourceForumPostId} are named in \texttt{library.md}; only owner-visible data is shown in the storage table.
\item \textbf{API and codes:} Library storage endpoint, statuses, and error codes are Not specified in proposals. \texttt{API.md} has no standalone library storage endpoint, and \texttt{response\_code.md} defines no library-specific error code.
\item \textbf{Pagination and limits:} Row limits, pagination, and refresh interval are Not specified in proposals.
\item \textbf{Notification and privacy rules:} No story-specific notification is stated. Storage access must respect owner privacy settings (\texttt{FR-27.5}, \texttt{NFR-PRIV-5}) and server-side authorization (\texttt{NFR-SEC-6}).
\end{itemize}
\par\medskip

\hypertarget{sec:b11-source-alignment-notes}{}
\addcontentsline{toc}{subsubsection}{B11. Source Alignment Notes}
\subsubsection*{B11. Source Alignment Notes}
\par\medskip

\begin{itemize}
\item This use case is derived from US-28 in \texttt{userstory.md} section 8.4 and from \texttt{FR-28.1} to \texttt{FR-28.3} and \texttt{AC-28.1} to \texttt{AC-28.3} in the proposals.
\item \texttt{library.md} supplies recommended component names, status values, and exact UI messages; these are UI recommendations, not new FRs.
\item \texttt{API.md} has no standalone library storage endpoint, so no storage route, status code, or error code is invented.
\item Overlap with existing \texttt{UC-06-share-learning-resource.md}: UC-06 publishes a resource and directs the owner to the storage view; UC-24 defines what that storage view displays. Deletion is covered by US-29 rather than this use case.
\item No I-xx, E-xx, or G-xx model relationship ID is introduced; flow headings such as A2 and E2 are specification flow labels only.
\end{itemize}
\par\medskip

\pagebreak
\hypertarget{sec:uc-25-delete-my-published-resource}{}
\addcontentsline{toc}{section}{UC-25 --- Delete My Published Resource}
\section*{UC-25 --- Delete My Published Resource}
\par\medskip

\textbf{Version:} \texttt{1.0}
\par\medskip

\textbf{Performed by:} [Nguyen Le Quoc Huy] | \textbf{Edited by:} [Nguyen Le Quoc Huy]
\par\medskip

\begingroup
\small
\renewcommand{\baselinestretch}{1.0}\selectfont
\renewcommand{\arraystretch}{1.15}
\setlength{\tabcolsep}{6pt}
\begin{longtable}{|>{\raggedright\arraybackslash}p{0.09\textwidth}|>{\raggedright\arraybackslash}p{0.84\textwidth}|}
\hline
\textbf{Field} & \textbf{Value} \\ \hline
\endhead
Use Case ID & UC-25 \\ \hline
Name & Delete My Published Resource \\ \hline
Primary Actor & University Student \\ \hline
Supporting Actors & None \\ \hline
Feature Group & Learning Resource Library \\ \hline
Priority & Not explicitly stated \\ \hline
Related User Stories & US-29 \\ \hline
Related Requirements & FR-29.1 to FR-29.4; AC-29.1 to AC-29.4; FR-28.1; FR-28.3; NFR-PERF-1; NFR-SCAL-4; NFR-SEC-1; NFR-SEC-6; NFR-PRIV-5; NFR-PRIV-7 \\ \hline
Source Documents & \texttt{docs/requirements/proposals/userstory.md}; \texttt{docs/requirements/proposals/functional\_requirement.md}; \texttt{docs/requirements/proposals/library.md}; \texttt{docs/requirements/proposals/create\_post.md}; \texttt{docs/requirements/proposals/API.md}; \texttt{docs/requirements/proposals/response\_code.md}; \texttt{docs/requirements/proposals/non\_functional\_requirement.md}; \texttt{docs/requirements/proposals/feature\_flow.md}; \texttt{docs/requirements/proposals/screen\_flow.md} \\ \hline
\end{longtable}
\endgroup
\par\medskip

\hypertarget{sec:part-a-use-case-model}{}
\addcontentsline{toc}{subsection}{Part A --- Use Case Model}
\subsection*{Part A --- Use Case Model}
\par\medskip

\hypertarget{sec:a1-focused-use-case-diagram}{}
\addcontentsline{toc}{subsubsection}{A1. Focused Use Case Diagram}
\subsubsection*{A1. Focused Use Case Diagram}
\par\medskip

\begin{center}
\includegraphics[width=0.85\textwidth,height=0.75\textheight,keepaspectratio]{"figures/UC-25-delete-published-resource_1.png"}
\end{center}
\par\medskip

\textbf{Diagram notation.} Person shape = human actor; cloud = supporting system; blue = selected use case; solid line = association; labelled solid arrow = generalization; dashed arrow = include or extend.
\par\medskip

\hypertarget{sec:a2-actor-register}{}
\addcontentsline{toc}{subsubsection}{A2. Actor Register}
\subsubsection*{A2. Actor Register}
\par\medskip

\begingroup
\small
\renewcommand{\baselinestretch}{1.0}\selectfont
\renewcommand{\arraystretch}{1.15}
\setlength{\tabcolsep}{6pt}
\begin{longtable}{|>{\raggedright\arraybackslash}p{0.11\textwidth}|>{\raggedright\arraybackslash}p{0.09\textwidth}|>{\raggedright\arraybackslash}p{0.71\textwidth}|}
\hline
\textbf{Actor} & \textbf{Type} & \textbf{Description} \\ \hline
\endhead
University Student & Primary & Authenticated resource owner who deletes a published learning resource that the student no longer wants to share. \\ \hline
\end{longtable}
\endgroup
\par\medskip

\hypertarget{sec:a3-relationships}{}
\addcontentsline{toc}{subsubsection}{A3. Relationships}
\subsubsection*{A3. Relationships}
\par\medskip

\hypertarget{sec:a3-1-associations}{}
\paragraph*{A3.1 Associations}
\par\medskip

\begingroup
\small
\renewcommand{\baselinestretch}{1.0}\selectfont
\renewcommand{\arraystretch}{1.15}
\setlength{\tabcolsep}{6pt}
\begin{longtable}{|>{\raggedright\arraybackslash}p{0.14\textwidth}|>{\raggedright\arraybackslash}p{0.26\textwidth}|>{\raggedright\arraybackslash}p{0.10\textwidth}|>{\raggedright\arraybackslash}p{0.42\textwidth}|}
\hline
\textbf{Actor} & \textbf{Use Case} & \textbf{Type} & \textbf{Description} \\ \hline
\endhead
University Student & UC-25 Delete My Published Resource & Primary & Confirms deletion of a resource the student published. \\ \hline
\end{longtable}
\endgroup
\par\medskip

\hypertarget{sec:a3-2-generalization-and-specialization}{}
\paragraph*{A3.2 Generalization and Specialization}
\par\medskip

\begingroup
\small
\renewcommand{\baselinestretch}{1.0}\selectfont
\renewcommand{\arraystretch}{1.15}
\setlength{\tabcolsep}{6pt}
\begin{longtable}{|>{\raggedright\arraybackslash}p{0.09\textwidth}|>{\raggedright\arraybackslash}p{0.09\textwidth}|>{\raggedright\arraybackslash}p{0.26\textwidth}|>{\raggedright\arraybackslash}p{0.47\textwidth}|}
\hline
\textbf{Child} & \textbf{Parent} & \textbf{Type} & \textbf{Inherited Behavior or Association} \\ \hline
\endhead
None & None & No generalization or specialization & The proposals name only the University Student actor for US-29. \\ \hline
\end{longtable}
\endgroup
\par\medskip

\hypertarget{sec:a3-3-include}{}
\paragraph*{A3.3 Include}
\par\medskip

\begingroup
\small
\renewcommand{\baselinestretch}{1.0}\selectfont
\renewcommand{\arraystretch}{1.15}
\setlength{\tabcolsep}{6pt}
\begin{longtable}{|>{\raggedright\arraybackslash}p{0.23\textwidth}|>{\raggedright\arraybackslash}p{0.30\textwidth}|>{\raggedright\arraybackslash}p{0.40\textwidth}|}
\hline
\textbf{Base Use Case} & \textbf{Included Use Case} & \textbf{Meaning} \\ \hline
\endhead
None & None & No include relationship & The proposals do not name an included use case for US-29. \\ \hline
\end{longtable}
\endgroup
\par\medskip

\hypertarget{sec:a3-4-extend}{}
\paragraph*{A3.4 Extend}
\par\medskip

\begingroup
\small
\renewcommand{\baselinestretch}{1.0}\selectfont
\renewcommand{\arraystretch}{1.15}
\setlength{\tabcolsep}{6pt}
\begin{longtable}{|>{\raggedright\arraybackslash}p{0.15\textwidth}|>{\raggedright\arraybackslash}p{0.11\textwidth}|>{\raggedright\arraybackslash}p{0.18\textwidth}|>{\raggedright\arraybackslash}p{0.47\textwidth}|}
\hline
\textbf{Extending Use Case} & \textbf{Base Use Case} & \textbf{Condition} & \textbf{Added Behavior} \\ \hline
\endhead
None & None & No extend relationship & The proposals do not name an extending use case for US-29. \\ \hline
\end{longtable}
\endgroup
\par\medskip

\hypertarget{sec:a4-model-notes}{}
\addcontentsline{toc}{subsubsection}{A4. Model Notes}
\subsubsection*{A4. Model Notes}
\par\medskip

\begin{enumerate}
\item The scope is owner-initiated deletion, confirmation, removal from search and library listings, and access restriction after deletion.
\item Publishing resources and viewing storage statistics are outside this use case; storage is US-28 and publishing is US-27.
\item The delete entry point is the owner's resource list or resource detail view, as recommended by \texttt{library.md}.
\item No distinct included, extending, or specialized use case is named for US-29, so no I-xx, E-xx, or G-xx relationship ID is introduced.
\item The exact API endpoint, route, deletion status code, and priority are not specified in proposals.
\end{enumerate}
\par\medskip

\hypertarget{sec:part-b-use-case-specification}{}
\addcontentsline{toc}{subsection}{Part B --- Use Case Specification}
\subsection*{Part B --- Use Case Specification}
\par\medskip

\hypertarget{sec:b1-brief-description}{}
\addcontentsline{toc}{subsubsection}{B1. Brief Description}
\subsubsection*{B1. Brief Description}
\par\medskip

A university student deletes a resource the student published so that content the student no longer wants to share is removed. The system requires confirmation before deletion, removes the deleted resource from search and library listings, and prevents other users from accessing it.
\par\medskip

\texttt{library.md} recommends placing the delete control in \textbf{My Storage}, with a dedicated confirmation dialog. The exact dialog text is specified in \texttt{library.md}; the functional rules come from \texttt{FR-29.1} to \texttt{FR-29.4} and \texttt{AC-29.1} to \texttt{AC-29.4}.
\par\medskip

\hypertarget{sec:b2-preconditions}{}
\addcontentsline{toc}{subsubsection}{B2. Preconditions}
\subsubsection*{B2. Preconditions}
\par\medskip

\begin{enumerate}
\item The student is authenticated before accessing protected user resources (\texttt{NFR-SEC-1}).
\item The Learning Resource Library is reachable through \texttt{Home -$>$ Library Group} in \texttt{screen\_flow.md}.
\item The student has a published resource available in the personal storage view (\texttt{FR-28.1}, \texttt{AC-28.1}).
\item The resource is owned by the student who requests deletion (\texttt{FR-29.1}, \texttt{AC-29.1}).
\item Search, library-listing, and access-enforcement mechanisms are available (\texttt{FR-29.3}, \texttt{FR-29.4}).
\end{enumerate}
\par\medskip

\hypertarget{sec:b3-postconditions}{}
\addcontentsline{toc}{subsubsection}{B3. Postconditions}
\subsubsection*{B3. Postconditions}
\par\medskip

\begin{enumerate}
\item \textbf{Success:} The confirmed resource is deleted from the owner's published resources.
\item \textbf{Success:} The deleted resource is removed from search and library listings.
\item \textbf{Success:} Other users cannot access the deleted resource.
\item \textbf{Failure:} No resource is deleted when confirmation is cancelled or authorization fails.
\item \textbf{Failure:} If deletion or removal fails, the resource is not presented as deleted and retry behavior is shown where supported.
\end{enumerate}
\par\medskip

\hypertarget{sec:b4-main-success-scenario-basic-flow}{}
\addcontentsline{toc}{subsubsection}{B4. Main Success Scenario (Basic Flow)}
\subsubsection*{B4. Main Success Scenario (Basic Flow)}
\par\medskip

\begin{enumerate}
\item The student opens \textbf{My Storage} through \texttt{MyStorageTab}.
\item The system displays \texttt{MyStoragePanel} with the student's published resources.
\item The student selects \textbf{Delete} through \texttt{DeleteMyResourceButton} for the chosen resource.
\item The system displays \texttt{DeleteResourceConfirmDialog} with the exact message “Delete \texttt{$<$title$>$}? It will be removed from search and library listings.”
\item The dialog provides \textbf{Cancel} and \textbf{Delete resource} buttons.
\item The student selects \textbf{Delete resource}.
\item The system verifies that the student owns the resource (\texttt{FR-29.1}, \texttt{AC-29.1}, \texttt{NFR-SEC-6}).
\item The system removes the resource from search and library listings (\texttt{FR-29.3}, \texttt{AC-29.3}).
\item The system prevents other users from accessing the deleted resource (\texttt{FR-29.4}, \texttt{AC-29.4}).
\item The system confirms deletion through \texttt{LibraryToastHost}; the exact success-message text is Not specified in proposals.
\end{enumerate}
\par\medskip

\hypertarget{sec:b5-alternative-flows}{}
\addcontentsline{toc}{subsubsection}{B5. Alternative Flows}
\subsubsection*{B5. Alternative Flows}
\par\medskip

\hypertarget{sec:a1-cancel-deletion}{}
\paragraph*{A1. Cancel deletion}
\par\medskip

The student selects \textbf{Cancel} in \texttt{DeleteResourceConfirmDialog}. The system closes the dialog and leaves the resource unchanged (\texttt{AC-29.2}, \texttt{FR-29.2}).
\par\medskip

\hypertarget{sec:a2-delete-from-resource-detail}{}
\paragraph*{A2. Delete from resource detail}
\par\medskip

The student opens the resource detail from \textbf{View} and starts deletion there. The same confirmation dialog and deletion rules apply.
\par\medskip

\hypertarget{sec:a3-return-to-storage-after-deletion}{}
\paragraph*{A3. Return to storage after deletion}
\par\medskip

After successful deletion, the student remains in or returns to \textbf{My Storage}, and the deleted resource no longer appears in the owned-resource list.
\par\medskip

\hypertarget{sec:b6-exception-flows}{}
\addcontentsline{toc}{subsubsection}{B6. Exception Flows}
\subsubsection*{B6. Exception Flows}
\par\medskip

\hypertarget{sec:e1-student-does-not-own-the-resource}{}
\paragraph*{E1. Student does not own the resource}
\par\medskip

The student attempts to delete a resource that is not owned by the student. The system rejects the deletion through authorization checks (\texttt{NFR-SEC-6}); the exact authorization message is Not specified in proposals.
\par\medskip

\hypertarget{sec:e2-deletion-or-listing-removal-fails}{}
\paragraph*{E2. Deletion or listing removal fails}
\par\medskip

The system cannot complete deletion or cannot remove the resource from search and library listings. The resource is not presented as deleted, and a retry path is shown where supported; the exact failure message is Not specified in proposals.
\par\medskip

\hypertarget{sec:e3-access-to-deleted-resource-attempted}{}
\paragraph*{E3. Access to deleted resource attempted}
\par\medskip

Another user attempts to access the deleted resource. The system blocks access (\texttt{FR-29.4}, \texttt{AC-29.4}) and may show the exact access-denied message “You do not have permission to access this resource.”
\par\medskip

\hypertarget{sec:b7-business-rules}{}
\addcontentsline{toc}{subsubsection}{B7. Business Rules}
\subsubsection*{B7. Business Rules}
\par\medskip

\begin{enumerate}
\item Users can delete their own published resources (\texttt{FR-29.1}, \texttt{AC-29.1}).
\item The system requires confirmation before deleting a published resource (\texttt{FR-29.2}, \texttt{AC-29.2}).
\item The deleted resource is removed from search and library listings (\texttt{FR-29.3}, \texttt{AC-29.3}).
\item Other users cannot access a deleted resource (\texttt{FR-29.4}, \texttt{AC-29.4}).
\item Publication status and date are owner-visible before deletion (\texttt{FR-28.3}, \texttt{AC-28.3}).
\item Privacy settings continue to govern shared resources until deletion is complete (\texttt{FR-27.5}, \texttt{NFR-PRIV-5}, \texttt{NFR-PRIV-7}).
\end{enumerate}
\par\medskip

\hypertarget{sec:b8-special-requirements}{}
\addcontentsline{toc}{subsubsection}{B8. Special Requirements}
\subsubsection*{B8. Special Requirements}
\par\medskip

\begingroup
\small
\renewcommand{\baselinestretch}{1.0}\selectfont
\renewcommand{\arraystretch}{1.15}
\setlength{\tabcolsep}{6pt}
\begin{longtable}{|>{\raggedright\arraybackslash}p{0.10\textwidth}|>{\raggedright\arraybackslash}p{0.83\textwidth}|}
\hline
\textbf{ID} & \textbf{Requirement} \\ \hline
\endhead
NFR-PERF-1 & The system shall provide responsive user interactions with a response time of ≤ 2 seconds for standard user actions. \\ \hline
NFR-SCAL-4 & The system architecture shall support at least 100,000 uploaded documents and learning resources. \\ \hline
NFR-SEC-1 & The system shall require authentication before accessing protected user resources. \\ \hline
NFR-SEC-6 & The system shall prevent unauthorized access through server-side authorization checks. \\ \hline
NFR-PRIV-5 & Users shall be able to control access to documents and learning resources they share. \\ \hline
NFR-PRIV-7 & Privacy settings shall apply consistently to forum posts, comments, member lists, and search results. \\ \hline
Design constraint & The delete control, confirmation dialog text, button labels, and post-deletion feedback follow \texttt{library.md} as UI recommendations; FR/AC/NFR entries are confirmed requirements. \\ \hline
\end{longtable}
\endgroup
\par\medskip

\hypertarget{sec:b9-extension-points}{}
\addcontentsline{toc}{subsubsection}{B9. Extension Points}
\subsubsection*{B9. Extension Points}
\par\medskip

\begingroup
\small
\renewcommand{\baselinestretch}{1.0}\selectfont
\renewcommand{\arraystretch}{1.15}
\setlength{\tabcolsep}{6pt}
\begin{longtable}{|>{\raggedright\arraybackslash}p{0.39\textwidth}|>{\raggedright\arraybackslash}p{0.54\textwidth}|}
\hline
\textbf{Extension Point} & \textbf{Location in Flow of Events} \\ \hline
\endhead
Cancel confirmation & Basic flow step 5; alternative A1 \\ \hline
Detail-page deletion & Before basic flow step 3; alternative A2 \\ \hline
Post-deletion storage refresh & After basic flow step 9; alternative A3 \\ \hline
Ownership failure & Basic flow step 7; exception E1 \\ \hline
Deletion or removal failure & Basic flow steps 8-9; exception E2 \\ \hline
Deleted-resource access & After basic flow step 9; exception E3 \\ \hline
\end{longtable}
\endgroup
\par\medskip

\hypertarget{sec:b10-interface-and-data-notes}{}
\addcontentsline{toc}{subsubsection}{B10. Interface and Data Notes}
\subsubsection*{B10. Interface and Data Notes}
\par\medskip

\begin{itemize}
\item \textbf{Screen and flow path (confirmed):} \texttt{screen\_flow.md} defines \texttt{Home -$>$ Library Group}; \texttt{userstory.md} section 8.5 defines the delete requirement. \texttt{library.md} flow 7.4 defines \texttt{My Storage -$>$ DeleteMyResourceButton -$>$ DeleteResourceConfirmDialog -$>$ removed from search and listings}.
\item \textbf{Route (not specified):} Route is Not specified in proposals; \texttt{library.md} presents \texttt{/library/my-storage} as a suggested state only.
\item \textbf{Recommended UI:} \texttt{MyStorageTab} label \textbf{My Storage}, \texttt{MyPublishedResourcesTable}, \texttt{DeleteMyResourceButton} label \textbf{Delete}, \texttt{MyResourceOpenButton} label \textbf{View}, and \texttt{DeleteResourceConfirmDialog}.
\item \textbf{Exact delete confirmation text:} “Delete \texttt{$<$title$>$}? It will be removed from search and library listings.” Buttons are \textbf{Cancel} and \textbf{Delete resource}.
\item \textbf{Other exact messages:} “You have not published any resources yet.” for empty storage; “Unable to load resources.” with \textbf{Retry} for storage load failure; “You do not have permission to access this resource.” for blocked access. The deletion success-message text is Not specified in proposals.
\item \textbf{Data fields:} \texttt{resourceId}, \texttt{title}, \texttt{resourceType}, \texttt{status}, \texttt{publishedAt}, \texttt{privacy}, \texttt{views}, and \texttt{accesses} are named in \texttt{library.md}. Status values are exactly \texttt{Published}, \texttt{Restricted}, and \texttt{Unavailable}.
\item \textbf{API and codes:} Library delete endpoint, statuses, and error codes are Not specified in proposals. \texttt{API.md} has no standalone library delete endpoint, and \texttt{response\_code.md} defines no library-specific error code; forum error codes must not be invented for library deletion.
\item \textbf{Authorization and privacy rules:} Only the owner may delete the resource (\texttt{FR-29.1}, \texttt{AC-29.1}, \texttt{NFR-SEC-6}); other users cannot access the deleted resource (\texttt{FR-29.4}, \texttt{AC-29.4}, \texttt{NFR-PRIV-5}, \texttt{NFR-PRIV-7}).
\end{itemize}
\par\medskip

\hypertarget{sec:b11-source-alignment-notes}{}
\addcontentsline{toc}{subsubsection}{B11. Source Alignment Notes}
\subsubsection*{B11. Source Alignment Notes}
\par\medskip

\begin{itemize}
\item This use case is derived from US-29 in \texttt{userstory.md} section 8.5 and from \texttt{FR-29.1} to \texttt{FR-29.4} and \texttt{AC-29.1} to \texttt{AC-29.4} in the proposals.
\item The confirmation dialog text, button labels, and storage entry point come from \texttt{library.md} as UI recommendations tied to \texttt{FR-29.2} and \texttt{AC-29.2}.
\item \texttt{API.md} has no standalone library delete endpoint, so no library delete route, status code, or error code is invented.
\item Overlap with existing \texttt{UC-06-share-learning-resource.md}: UC-06 establishes that published resources must become searchable and must be removed from listings and blocked from access when deleted; UC-25 defines the owner's complete deletion flow.
\item No I-xx, E-xx, or G-xx model relationship ID is introduced; flow headings such as A2 and E2 are specification flow labels only.
\end{itemize}
\par\medskip

\end{document}
